\providecommand{\recsinc}{\operatorname{sinc}}
\providecommand{\recgE}{\gamma_{\mathrm{E}}}
\providecommand{\recGen}{\operatorname{Gen}}
\providecommand{\recav}{\operatorname{av}}
\providecommand{\recA}{\mathcal{A}}
\providecommand{\recartanh}{\operatorname{artanh}}
\providecommand{\anapos}[1]{\bigl(#1\bigr)_{+}}
\providecommand{\ananeg}[1]{\bigl(#1\bigr)_{-}}
\providecommand{\anaone}{\mathbf{1}}
\providecommand{\anaE}{\mathcal{E}}
\providecommand{\anaEx}{\mathbb{E}}
\providecommand{\anaPr}{\mathbb{P}}
\providecommand{\anaM}{\mathcal{M}}
\providecommand{\anaW}{\mathcal{W}}
\providecommand{\anaZ}{\mathcal{Z}}
\providecommand{\anaC}{\mathcal{C}}
\providecommand{\liE}{\mathcal{E}}
\providecommand{\liD}{\mathcal{D}}
\providecommand{\liZ}{\mathcal{Z}}
\providecommand{\liT}{\mathbb{T}}
\providecommand{\lidist}[1]{\lVert #1\rVert_{\mathbb{T}}}
\providecommand{\liEx}{\mathbb{E}}
\providecommand{\liind}{\mathbf{1}}
\providecommand{\liHk}{\mathcal{H}}
\providecommand{\liKk}{\mathcal{K}}
\section{The Li heat semigroup and random laws built from zeros}\label{sec:lih}

\subsection{A universal process from the Cayley angles}\label{li:sec:univ}
\sources{rh_li_universality_2026-11-26}

In this subsection $\liZ$ is a locally finite multiset of real numbers with multiplicities, and
$N(T)=\#\{\gamma\in\liZ:\abs\gamma\le T\}$ satisfies, for constants $d>0$ and $B\in\R$,
\begin{equation}\label{li:eq:count}
 N(T)=\frac d\pi T\log T+\frac B\pi T+O(\log(T+2))\qquad(T\to\infty).
\end{equation}
Put $\theta(\gamma)=2\arctan(1/(2\gamma))$ for $\gamma\ne0$, $\theta(0)=\pi$,
$\nu=\frac12\sum_{\gamma\in\liZ}(\delta_{\theta(\gamma)}+\delta_{-\theta(\gamma)})$, so that $m$ occurrences
of $0$ give mass $m/2$ at each of $\pm\pi$, and
\begin{equation}\label{li:eq:symbol}
 \Psi(u)=\int_\R(1-\cos uy)\,\nu(\dd y)=\sum_{\gamma\in\liZ}(1-\cos(u\theta(\gamma))).
\end{equation}
Euler's constant is $\gamma_E$ here.

\begin{lemma}[The process; \stU]\label{li:lem:levy}
$\nu$ is a symmetric L\'evy measure on $[-\pi,\pi]$ with $\nu(\R)=\infty$, $\int\abs y\,\nu(\dd y)=\infty$
and $0<V:=\int y^2\,\nu(\dd y)<\infty$. There is a symmetric L\'evy process $(X_t)_{t\ge0}$ with
$\liEx e^{iuX_t}=e^{-t\Psi(u)}$; it has no Gaussian part, infinitely many jumps, infinite variation,
and $\liEx X_t^2=tV$.
\end{lemma}

\begin{proof}
By \eqref{li:eq:count} and partial summation $\sum_{\abs\gamma>1}\gamma^{-2}<\infty$ and
$\sum_{\abs\gamma>1}\abs\gamma^{-1}=\infty$, and $\abs{\theta(\gamma)}\sim\abs\gamma^{-1}$. L\'evy--Khintchine
gives the process with zero Gaussian coefficient and drift, the variation criterion gives infinite
variation, and two derivatives of the characteristic function at zero give the variance.
\end{proof}

\begin{lemma}[The symbol at large frequency; \stU]\label{li:lem:symbol}
With $A=d/2$ and $\kappa=\gamma_E+B/d$,
\[
 \Psi(u)=A\abs u(\log\abs u+\kappa)+O(\sqrt{\abs u}\log(\abs u+2))\qquad(\abs u\to\infty).
\]
\end{lemma}

\begin{proof}
Take $u\ge2$ and $T=\sqrt u$. Ordinates with $\abs\gamma\le T$ contribute $O(T\log(T+2))$. Beyond,
$\theta(\gamma)=\gamma^{-1}+O(\abs\gamma^{-3})$, and replacing $\theta$ by $1/\gamma$ costs
$O(u\int_T^\infty v^{-3}\,\dd N(v))=O(\log u)$. Let $M(v)=(dv\log v+Bv)/\pi$ and $R=N-M$. For
$g(v)=1-\cos(u/v)$, integration by parts and $\abs{g'(v)}\le u/v^2$ give
$\int_T^\infty g\,\dd R=O(\log(T+2)+u\int_T^\infty\log(v+2)v^{-2}\,\dd v)=O(\sqrt u\log u)$, and extending the
main integral down to zero costs $O(T\log(T+2))$. Hence
\begin{align*}
 \Psi(u)&=\frac1\pi\int_0^\infty(1-\cos(u/v))(d\log v+d+B)\,\dd v+O(\sqrt u\log u)\\
 &=\frac u\pi[(d\log u+d+B)I_0-dI_1]+O(\sqrt u\log u),
\end{align*}
with $I_0=\int_0^\infty(1-\cos y)y^{-2}\dd y=\pi/2$ and
$I_1=\int_0^\infty(1-\cos y)y^{-2}\log y\,\dd y=\frac\pi2(1-\gamma_E)$. Both follow by differentiating
$\int_0^\infty(1-\cos y)y^{-1-\alpha}\dd y=\pi/(2\Gamma(1+\alpha)\sin(\pi\alpha/2))$, $0<\alpha<2$, at
$\alpha=1$; this identity follows from the Laplace representation of $y^{-1-\alpha}$, Tonelli's theorem
and $\int_0^\infty e^{-ty}(1-\cos y)\dd y=1/[t(1+t^2)]$.
\end{proof}

Let $W_0$ be the positive branch of Lambert's function and put
\begin{equation}\label{li:eq:scale}
 w(t)=W_0\Bigl(\frac{e^\kappa}{At}\Bigr),\qquad a(t)=Atw(t),\qquad\text{so}\quad a=e^{\kappa-w},\quad
 \log(1/a)+\kappa=w,\quad t=\frac a{Aw}.
\end{equation}
For $k\ge0$ define, with the principal logarithm,
\begin{equation}\label{li:eq:profiles}
 Q_k(x)=\frac{(-1)^k}{\pi k!}\Re\Bigl.\frac{\partial^k}{\partial z^k}\bigl[\Gamma(z)(1-ix)^{-z}\bigr]\Bigr|_{z=k+1}.
\end{equation}

\begin{theorem}[Universal expansion at every logarithmic order; \stU]\label{li:thm:allorders}
For every $t>0$, $X_t$ has a density $f_t$ that extends to an entire function. For fixed integers
$K,r\ge0$, as $t\downarrow0$,
\[
 \sup_{x\in\R}\Bigl|\frac{\dd^r}{\dd x^r}\Bigl[a(t)f_t(a(t)x)-\sum_{k=0}^K\frac{Q_k(x)}{w(t)^k}\Bigr]\Bigr|
 =O_{K,r}\bigl(w(t)^{-K-1}+\sqrt{a(t)}\bigr)=O_{K,r}(w(t)^{-K-1}).
\]
The constants depend on the multiset and its counting estimate.
\end{theorem}

\begin{proof}
\Cref{li:lem:symbol} makes $\exp(-t\Psi(u)+c\abs u)$ integrable for all $c>0$, so Fourier inversion
gives an entire density. Put $F_w(s)=s+s\log s/w$ for $s>0$; by \eqref{li:eq:scale},
$tA\frac sa(\log\frac sa+\kappa)=F_w(s)$. We first show
\begin{equation}\label{li:eq:fcomp}
 \int_0^\infty s^r\bigl|e^{-t\Psi(s/a)}-e^{-F_w(s)}\bigr|\,\dd s=O_r(\sqrt a).
\end{equation}
On $0<s<\sqrt a$ both exponentials are bounded (the negative part of $F_w$ is at most $e^{-w-1}/w$).
On $s\ge\sqrt a$ both exponents are at least $cs$ for small $t$, since $F_w(s)/s\ge\frac12+\kappa/(2w)$
and the remainder in \cref{li:lem:symbol} is a small fraction of its main term for
$s/a\ge a^{-1/2}$; the exponential mean value bound gives
$\abs{e^{-t\Psi(s/a)}-e^{-F_w(s)}}\le Ct\sqrt{s/a}\log(2+s/a)e^{-cs}$, and $t=a/(Aw)$ with
$\log(2+s/a)\le w+C_\kappa+\log(2+s)$ gives \eqref{li:eq:fcomp}. Next, Taylor's remainder for
$e^{-z}$, $z\ge0$, and $\abs{s\log s}\le1/e$ on $(0,1)$ give
\[
 \int_0^\infty s^re^{-s}\Bigl|e^{-s\log s/w}-\sum_{k=0}^K\frac{(-s\log s)^k}{k!\,w^k}\Bigr|\,\dd s=O_{K,r}(w^{-K-1}).
\]
These are $L^1$ bounds for the Fourier transforms with polynomial weights; Fourier inversion and
differentiation give the uniform claim, and
$\int_0^\infty s^k(\log s)^ke^{-(1-ix)s}\dd s=\partial_z^k[\Gamma(z)(1-ix)^{-z}]_{z=k+1}$ identifies the
profiles. As $a=e^{\kappa-w}$, $\sqrt a$ is smaller than every power of $1/w$.
\end{proof}

The first profiles are
\begin{gather*}
 Q_0(x)=\frac1{\pi(1+x^2)},\qquad Q_1(x)=-\frac1\pi\Re\frac{1-\gamma_E-\log(1-ix)}{(1-ix)^2},\\
 Q_2(x)=\frac1\pi\Re\frac{(\frac32-\gamma_E-\log(1-ix))^2+\frac{\pi^2}6-\frac54}{(1-ix)^3},
\end{gather*}
so the leading shape is the standard Cauchy density.

\begin{corollary}[Universality beyond every logarithmic order; \stU]\label{li:cor:universal}
If two multisets satisfy \eqref{li:eq:count} with the same $d,B$, their scaled densities differ by
$O_r(\sqrt a)$ in $C_b^r(\R)$, and in particular have the same expansion to every order in
$1/\log(1/t)$.
\end{corollary}

\begin{proof}
Both scaled characteristic functions differ from $e^{-F_w(\abs s)}$ by $O_r(\sqrt a)$ in the weighted
$L^1$ norm \eqref{li:eq:fcomp}; apply Fourier inversion.
\end{proof}

For the ordinates of a fixed primitive Dirichlet $L$-function of conductor $q$, counted with
multiplicity and both signs, $N_\chi(T)=\frac T\pi\log\frac{qT}{2\pi e}+O_\chi(\log(T+2))$
\cite{BennettMartinOBryantRechnitzer2021Counting}, and $q=1$ gives $\zeta$. Thus $d=1$,
$B=\log(q/(2\pi))-1$, and
\[
 w_q(t)=W_0\Bigl(\frac{qe^{\gamma_E-1}}{\pi t}\Bigr),\qquad a_q(t)=\frac t2w_q(t)\sim\tfrac12t\log(1/t).
\]
These density statements are unconditional, and all primitive characters of one conductor share the
expansion to every logarithmic order; no uniformity in $q$ is asserted. By \cref{src:thm:dgrh}
every zero is $\rho=\frac12+i\gamma$ with $1-1/\rho=e^{i\theta(\gamma)}$ (with $e^{i\pi}=-1$ at a central zero),
so the Li coefficients satisfy
\begin{equation}\label{li:eq:liident}
 \Re\lambda_\chi(n)=\sum_\rho\Re\bigl[1-(1-1/\rho)^n\bigr]=\Psi_\chi(n),\qquad
 \liEx e^{inX_t}=e^{-t\Re\lambda_\chi(n)}\quad(\text{\stG}),
\end{equation}
and $\liEx e^{inX_t}=e^{-t\lambda_n}$ for $\zeta$. For nonreal characters the real part is essential.

Reduce $X_t$ modulo $2\pi$, with density $p_t(y)=\frac1{2\pi}\sum_{n\in\Z}e^{-t\Psi(n)}e^{iny}$ on
$[-\pi,\pi]$, an entire periodic function. By \eqref{li:eq:liident} the generator $\mathcal A$ of this
circle semigroup has eigenvalues $\Re\lambda_\chi(\abs n)$.

\begin{proposition}[Circle kernel and heat trace; \stU]\label{li:prop:trace}
For fixed $R>0$ and $K,r\ge0$,
$\sup_{\abs x\le R}\bigl|\frac{\dd^r}{\dd x^r}[ap_t(ax)-\sum_{k\le K}Q_k(x)w^{-k}]\bigr|=O_{R,K,r}(w^{-K-1}+\sqrt a)$,
and
\[
 \Tr(e^{-t\mathcal A})=2\pi p_t(0)=\frac2a\Bigl[1-\frac{1-\gamma_E}w+\frac{(\frac32-\gamma_E)^2+\frac{\pi^2}6-\frac54}{w^2}
 +O(w^{-3})\Bigr].
\]
For a primitive Dirichlet character the leading term is $4/(tw_q(t))$.
\end{proposition}

\begin{proof}
The scaled series has frequency mesh $a$. Splitting at $an=\sqrt a$ and the estimates of
\eqref{li:eq:fcomp} bound its difference from the model Riemann sum by $O_r(\sqrt a)$. The model
function $s^re^{-F_w(s)}e^{ixs}$ has uniformly integrable derivative for $\abs x\le R$, so the
Riemann-sum error is $O_{R,r}(a)$; the Taylor step is unchanged, and $x=0$ gives the trace.
\end{proof}

Uniformity over all real $x$ fails on the circle: periodicity produces returning peaks at distance
$2\pi/a$.

\begin{theorem}[Moment transition; \stU]\label{li:thm:moments}
Let $L=\log(1/t)$. As $t\downarrow0$,
\[
 \liEx\abs{X_t}^p\sim\frac{a(t)^p}{\cos(\pi p/2)}\ (0<p<1),\qquad
 \liEx\abs{X_t}\sim\frac d{2\pi}tL^2,\qquad
 \liEx\abs{X_t}^p\sim t\int\abs y^p\,\nu(\dd y)\ (1<p<2),
\]
and $\liEx X_t^2=tV$. Moreover $X_t/a(t)\Rightarrow\mathrm{Cauchy}(0,1)$ as $t\downarrow0$ and
$X_t/\sqrt{tV}\Rightarrow N(0,1)$ as $t\to\infty$.
\end{theorem}

\begin{proof}
For $0<p<2$, $\liEx\abs{X_t}^p=c_p\int_0^\infty(1-e^{-t\Psi(u)})u^{-p-1}\dd u$ with
$c_p=2\Gamma(1+p)\sin(\pi p/2)/\pi$, by Tonelli and the cosine integral of \cref{li:lem:symbol}. For
$0<p<1$ put $u=s/a$: the range $s\le aU_0$ contributes $O(ta^{-p})=o(1)$ to the normalized moment by
$\Psi(u)\le Vu^2/2$; on $aU_0\le s\le1$, $t\Psi(s/a)\le Cs$, and on $s\ge1$, $1-e^{-t\Psi}\le1$;
dominated convergence gives the limit $c_p\int_0^\infty(1-e^{-s})s^{-p-1}\dd s=1/\cos(\pi p/2)$. For
$1<p<2$ divide by $t$ and use $(1-e^{-t\Psi})/t\le\Psi$ with the majorant from $\Psi=O(u^2)$ at $0$ and
\cref{li:lem:symbol} at infinity. For $p=1$ the range $u<1$ is $O(t)$; splitting at $1/t$ with
$1-e^{-z}\le\min(z,1)$ gives the upper bound $\frac A\pi tL^2+O(tL)$, and keeping
$1\le u\le1/(tL^2)$ with $1-e^{-z}\ge z-z^2/2$ gives the lower bound $\frac A\pi t(L-2\log L)^2+O(tL)$.
The weak limits follow from $t\Psi(s/a)\to\abs s$ and from $\Psi(u)=Vu^2/2+o(u^2)$ at $0$.
\end{proof}

The Cauchy limit has no variance because second moments are not uniformly integrable under the
small-time scaling. Nakamura and Suzuki relate RH to infinitely divisible distributions built from
ordinates with inverse-square weights \cite{NakamuraSuzuki2023Infinitely}; here the jumps are the
transformed angles with equal weights. General small-time moment asymptotics for L\'evy processes are
in \cite{FigueroaLopez2008Small}.

The densities of \cref{li:cor:universal} are one-time laws seen at the central scale. A single
exact time slice of the circle kernel already determines the character up to conjugation
(\cref{rec:thm:liheat}); \cref{li:sec:path} shows that one path does so almost surely, although no
finite set of times does, and \cref{li:sec:offorigin} describes the oriented kernel away from the
origin.




\subsection{The character from one Li heat kernel}\label{rec:sec:liheat}
\sources{dirichlet_li_heat_kernel; dirichlet_li_recovery}

Exponentiating the complex Li coefficients gives a positive convolution semigroup on the circle
$\mathbb T=\R/2\pi\Z$ whose jumps are the Cayley angles of the zeros; \cref{li:sec:univ} treats the
symmetrized semigroup with exponent $\Re\lambda_\chi$. Here the orientation is kept. Let $\chi$ be primitive of
conductor $q$, $q=1$ included. Under \cref{src:thm:dgrh} each nontrivial zero is $\rho=\tfrac12+i\gamma$ and
\[
 1-\rho^{-1}=e^{i\theta_\gamma},\qquad \theta_\gamma=2\arctan\frac1{2\gamma},\qquad\gamma=\tfrac12\cot(\theta_\gamma/2),
\]
with $\theta_0=\pi$; the point $0\in\mathbb T$ corresponds to no zero. Put, for $n\in\Z$ and $t>0$,
\[
 \lambda_\chi(n)=\lim_{T\to\infty}\sum_{\abs\gamma\le T}\bigl(1-e^{in\theta_\gamma}\bigr),\qquad
 p_{\chi,t}(y)=\frac1{2\pi}\sum_{n\in\Z}e^{-t\lambda_\chi(n)-iny},
\]
zeros repeated with multiplicity, so that $\lambda_\chi(-n)=\overline{\lambda_\chi(n)}$; the symmetric truncation matters
for the imaginary part, and the normalization through logarithmic derivatives at $s=1$ reverses the
index for complex $\chi$ \cite{Lagarias2007Li}. Let $\nu_\chi=\sum_\gamma\delta_{\theta_\gamma}$ on $(-\pi,\pi]$. The one-sided counts
$N^\pm_\chi(T)=\frac T{2\pi}\log\frac{qT}{2\pi e}+O_\chi(\log(2+T))$ give $\int\theta^2\dd\nu_\chi<\infty$, and Lagarias's asymptotic
under GRH is
\begin{equation}\label{rec:eq:lagarias}
 \lambda_\chi(n)=\tfrac12n\log n+C_\chi n+O_\chi\bigl(\sqrt n\log(2+n)\bigr)\qquad(n\to+\infty),\quad C_\chi\in\R
\end{equation}
\cite{Lagarias2007Li}. Let $w=W_0(qe^{\recgE-1}/(\pi t))$ with $W_0$ the positive branch of Lambert's function, and
$a=a(t)=\tfrac12tw\sim\tfrac12t\log(1/t)$.

\begin{proposition}[The oriented semigroup; \stG]\label{rec:prop:liheat}
$p_{\chi,t}$ is an entire periodic function, strictly positive on $\mathbb T$, of integral one, with
$\int_{\mathbb T}p_{\chi,t}(y)e^{iny}\dd y=e^{-t\lambda_\chi(n)}$; the laws form a convolution semigroup with jump measure $\nu_\chi$. As
$t\downarrow0$, $ap_{\chi,t}(ax)\to1/(\pi(1+x^2))$ locally uniformly and $p_{\chi,t}(0)\sim1/(\pi a)$.
\end{proposition}

\begin{proof}
For a finite cutoff the sum of independent unit-rate Poisson jumps at the listed angles has the
truncated exponent; the pointwise limit is positive definite and equals one at $n=0$, so it is the
Fourier sequence of a probability on $\mathbb T$, with compensated exponent
$\int(1-e^{iu\theta}+iu\theta)\dd\nu_\chi(\theta)-iub$, $b=\lim_T\sum_{\abs\gamma\le T}\theta_\gamma$, which converges because
$\sum^*\sin\theta_\gamma=-\Im\lambda_\chi(1)$ and $\sum\abs{\theta_\gamma-\sin\theta_\gamma}<\infty$. By \eqref{rec:eq:lagarias} the series converges
absolutely in every horizontal strip. A nonzero real-analytic nonnegative density is positive almost
everywhere, so $p_t=p_{t/2}*p_{t/2}$ is positive everywhere. The central profile for $\Re\lambda_\chi$ is
\cref{li:prop:trace}; replacing $\Re\lambda_\chi$ by $\lambda_\chi$ changes the scaled density by at most
$\frac{at}{2\pi}\sum_n\abs{\Im\lambda_\chi(n)}e^{-t\Re\lambda_\chi(n)}=O_\chi(\sqrt a)$, by \eqref{rec:eq:lagarias}: on $\abs n\ge a^{-1/2}$ the
exponential is at most $e^{-ca\abs n}$, and the smaller modes give $O(\sqrt a)$ directly.
\end{proof}

Write $F_{\chi,t}(n)=e^{-t\lambda_\chi(n)}$. A sequence converges in natural density to $u$ if for each $\eps>0$ the
proportion of $1\le n\le N$ with $\abs{u_n-u}>\eps$ tends to zero.

\paragraph{Finite blocks on a progression.}
Several characters, unknown positive or complex weights and an arithmetic progression of indices are
treated at once. Write $z_\rho=1-1/\rho=e^{i\theta_\gamma}$ and $m_\chi(\rho)$ for the multiplicity, sums over $\rho$ running
over distinct zeros. The count \eqref{rec:eq:licount} gives $\sum_\rho m_\chi(\rho)\abs{1-z_\rho}^2<\infty$, the starred first sum
converges \cite{Lagarias2007Li}, and the only possible limit point of the $z_\rho$ on the circle is $1$. Fix a nonempty
finite set $\mathcal C$ of distinct primitive characters, a stride $h\ge1$, a residue $0\le r<h$ and an integer $d\ge0$, and
for $b\in\C^{\mathcal C}$ put
\[
 A_b(k)=\sum_{\chi\in\mathcal C}b_\chi\lambda_\chi(r+hk)\qquad(k\in\Z).
\]
For positive weights $\exp(-A_b)$ is a product of heat coefficients $F_{\chi,b_\chi}(r+hk)$; complex weights allow
comparison and subtraction without choosing a branch of a complex power. A calibration polynomial $P$
multiplies the observed sequence by $\exp(P(k))$; it includes an unknown normalization and rotation. Let
$\Delta x(k)=x(k+1)-x(k)$.

\begin{theorem}[Uniform separation on finite blocks; \stG]\label{rec:thm:liblock}
Let $M>0$.
\begin{enumerate}[label=\textup{(\roman*)}]
\item There are an integer $W\ge1$ and $c>0$, depending only on $\mathcal C,h,r,d,M$, such that every interval $I$ of $W$
consecutive positive integers and every $b$ with $\norm b_2\le M$ satisfy
\begin{equation}\label{rec:eq:calibration}
 \inf_{\substack{P\in\C[X]\\ \deg P\le d}}\ \max_{k\in I}\bigl\lvert\exp\bigl(-A_b(k)+P(k)\bigr)-1\bigr\rvert\ \ge\ c\norm b_2 .
\end{equation}
The same constants apply when a different calibration polynomial is fitted on each interval.
\item With $j=\max(2,d+1)$ there are $L\ge1$ and $\eta_0>0$ such that for every $b$ with $0<\norm b_2\le M$ every block of
$L$ consecutive integers contains $k$ with $\abs{\exp(-\Delta^jA_b(k))-1}\ge\eta_0\norm b_2$.
\end{enumerate}
\end{theorem}

For a catalogue of $\abs{\mathcal C}$ characters every admissible $W$ satisfies $W\ge\abs{\mathcal C}+d+1$: on $W$ points the
values of the polynomials of degree at most $d$ form a subspace of dimension $\min(d+1,W)$; if its codimension is
smaller than $\abs{\mathcal C}$, the linear map $b\mapsto A_b|_I$ followed by the quotient has a nonzero kernel element,
which can be scaled to norm at most $M$ and whose values can be interpolated to make the left side of
\eqref{rec:eq:calibration} zero.

\begin{proof}
The sequence $u_b=\Delta^jA_b$ has the absolutely convergent expansion
\begin{equation}\label{rec:eq:pushprog}
 u_b(k)=-\sum_zB_zz^r(z^h-1)^j(z^h)^k,\qquad B_{z_\rho}=\sum_{\chi\in\mathcal C}b_\chi m_\chi(\rho),
\end{equation}
since $\abs{z^h-1}\le h\abs{1-z}$ and $j\ge2$. So $u_b$ is uniformly almost periodic on $\Z$ and depends continuously
on $b$ in the supremum norm. Suppose $\exp(-u_b(k))=1$ for all $k\in\Z$. Then $u_b$ takes values in $2\pi i\Z$, and an
almost period with error less than $2\pi$ is an exact period $v\ge1$. By uniqueness of Fourier coefficients of
finite measures on the circle, the measure in \eqref{rec:eq:pushprog}, grouped at $w=z^h$, is supported on $w^v=1$.
Choose an open arc about $1$ on which $z\mapsto z^h$ is injective; finitely many atoms lie outside it. Removing
their images and the $v$th roots of unity, above every remaining $w$ a contributing atom is the only preimage in the
arc, its coefficient in \eqref{rec:eq:pushprog} vanishes and $w\ne1$, so $B_z=0$. The discarded set is finite with finite
fibres, so $B_z=0$ except at finitely many $z$, and \cref{rec:lem:char} gives $b=0$. In particular $u_b\ne0$ when $b\ne0$.

For $\norm a_2=1$ and $0\le s\le M$ put $g_{a,s}=(\exp(-su_a)-1)/s$ for $s>0$ and $g_{a,0}=-u_a$. These are uniformly almost
periodic, $(a,s)\mapsto g_{a,s}$ is continuous into the bounded sequences with the supremum norm (at $s=0$ by the
exponential series), and each is nonzero by the preceding paragraph. For a nonzero uniformly almost periodic
sequence choose an index where its modulus exceeds $2\eta>0$; its translates by $\eta$-almost periods have bounded
gaps, so an excursion above $\eta$ occurs in every interval of a fixed length, and the same length works with
threshold $\eta/2$ throughout a supremum-norm neighbourhood. A finite subcover of the compact parameter space gives
$L$ and $\eta_0$ as in (ii), since $\exp(-u_b)-1=sg_{a,s}$ for $b=sa$.

For (i) write $R(k)=\exp(-A_b(k)+P(k))$. Since $\Delta^jP=0$,
\[
 \prod_{i=0}^jR(k+i)^{(-1)^{j-i}\binom ji}=\exp(-u_b(k)).
\]
This rational function of $j+1$ inputs equals $1$ at
$(1,\dots,1)$, and on the closed polydisc $\abs{R_i-1}\le\frac12$ its first derivatives are bounded, so its distance from
$1$ is at most $C_j\max_i\abs{R_i-1}$. Take $W=L+j$ and choose $k$ among the first $L$ positions of $I$ by (ii). Either some
$\abs{R(k+i)-1}>\frac12$, or $\max_i\abs{R(k+i)-1}\ge(\eta_0/C_j)\norm b_2$. Take $c=\min(1/(2M),\eta_0/C_j)$.
\end{proof}

Suppose two weighted lists from the fixed catalogue give sequences whose ratio, after calibration, differs from
$1$ by at most $\eps$ on a block of length $W$. If their weight difference has norm at most $M$, it has norm at most
$\eps/c$; at zero error the weights agree. At a common known time this recovers every factor multiplicity of a
primitive product. If the time is unknown, the identifiable quantities are time times multiplicity: $P$ at time
$at$ and $P^a$ at time $t$ give the same data for every positive integer $a$. For two individual finite mixtures
take the union of their catalogues and choose $M$ after the weights are given. Consequently a sequence of blocks
whose union has natural density zero determines a finite mixture, provided their lengths tend to infinity and
the calibrated ratio tends to $1$ in relative measure on them with calibrations of fixed degree: for distinct
mixtures every full subblock of length $W$ contains an error of at least $c\norm b_2$, a positive lower fraction. For
example the modes $r+hk$ with $2^{i^2}\le k<2^{i^2}+i$, $i\ge1$, suffice in this asymptotic sense. These statements control
relative errors of the Fourier coefficients, which may themselves be very small. Banks compares zero divisors
of two distinct primitive characters of one conductor directly, through short height windows
\cite{Banks2023Comparing}.

\begin{corollary}[One time slice determines the character; \stG]\label{rec:thm:liheat}
Let $\chi,\psi$ be primitive, $s,t>0$ and $\alpha\in\mathbb T$.
\begin{enumerate}
\item If $F_{\chi,t}(n)/(e^{in\alpha}F_{\psi,s}(n))\to1$ in natural density on the positive integers, then $s=t$, $\chi=\psi$ and $\alpha=0$.
\item If $\abs{F_{\chi,t}(n)}/\abs{F_{\psi,s}(n)}\to1$ in natural density, then $s=t$ and $\psi\in\set{\chi,\bar\chi}$, and both possibilities
give the same moduli.
\end{enumerate}
Hence one exact time slice $p_{\chi,t}$, even restricted to a nonempty open arc and known only up to
rotation, determines $t$ and $\chi$; if the orientation of the circle is also unknown, the ambiguity is
exactly conjugation.
\end{corollary}

\begin{proof}
Apply \cref{rec:thm:liblock}(i) with $h=1$, $r=0$, $d=1$ and $M=\norm b_2$. In (i) take $\mathcal C=\set{\chi,\psi}$ and
$b=(t,-s)$ if $\chi\ne\psi$, and $\mathcal C=\set\chi$, $b=t-s$ otherwise, and the calibration $P(k)=-ik\alpha$ for a real
representative $\alpha$. If $b\ne0$, every interval of $W$ consecutive positive integers contains an index with
$\abs{F_{\chi,t}/(e^{in\alpha}F_{\psi,s})-1}\ge c\norm b_2$, so these indices have lower density at least $1/W$, which contradicts
the hypothesis. Hence $\chi=\psi$ and $s=t$, so $e^{-in\alpha}\to1$ in density, its Ces\`aro mean is one, and $e^{i\alpha}=1$.
In (ii), conjugation reflects the ordinates, so $\lambda_{\bar\chi}(n)=\overline{\lambda_\chi(n)}$ and
$\abs{F_{\chi,t}(n)}=\exp(-\frac t2(\lambda_\chi+\lambda_{\bar\chi})(n))$. The quotient of the moduli is $\exp(-A_b)$, with $P=0$, for
$\mathcal C$ the set of distinct characters among $\chi,\bar\chi,\psi,\bar\psi$ and $b_\xi$ equal to $t/2$ times the number of
occurrences of $\xi$ in the list $\chi,\bar\chi$ minus $s/2$ times its number in the list $\psi,\bar\psi$. Two conjugacy orbits
are equal or disjoint; if they are disjoint $b\ne0$, and if they are equal $b=0$ only for $s=t$. The same density
argument concludes, and conjugation gives the same moduli. Equality on an
open arc extends to $\mathbb T$ by analyticity, and reflection of $y$ sends $p_{\chi,t}$ to $p_{\bar\chi,t}$.
\end{proof}

The relative form of the hypothesis is essential, since the high Fourier coefficients are extremely
small; and a general circle semigroup can lose its generator at one time: unit jumps of rate $\pi/t$ at
$\pi/2$, or at $-\pi/2$, have time-$t$ coefficients with ratio $\exp(2\pi i\sin(n\pi/2))=1$ \cite{HasebeHotta2022Additive}.
\Cref{li:thm:times} describes all such losses for integral jump measures, and \cref{rec:lem:char} removes exactly
this obstruction for Dirichlet zeros. By \cref{rec:thm:liblock}(ii) with $d=0$, applied to the same weights, the bounded
rotation-invariant sequence $F_{\chi,t}(n+1)F_{\chi,t}(n-1)/F_{\chi,t}(n)^2$ alone determines $t$ and $\chi$, and its modulus $t$
and the pair $\set{\chi,\bar\chi}$, also when two such observations agree only in density (\stG).

\begin{theorem}[An explicit block for $\zeta$ and Dirichlet beta; \stG]\label{rec:thm:li255}
Let $\beta(s)=L(s,\chi_{-4})$; the coefficients of $\zeta$ and $\beta$ are real, so their Li coefficients are real. For
$0\le r<20$ write $A_{u,v}(k)=u\lambda_\zeta(r+20k)+v\lambda_\beta(r+20k)$. For every interval $I$ of $255$ consecutive positive
integers, all real $u,v$ and every real affine function $P$,
\[
 \max_{k\in I}\abs{A_{u,v}(k)-P(k)}\ \ge\ \tfrac16\max(\abs u,\abs v).
\]
Consequently the logarithms of the heat-coefficient moduli of $\zeta^a\beta^b$ at time one, at these $255$ indices in any
residue class modulo $20$, determine the nonnegative integers $a,b$ even in the presence of an unknown affine
drift, and uniform errors strictly below $1/12$ in the logarithms permit at most one such pair.
\end{theorem}

\begin{proof}
The first two positive ordinates of $\zeta$ lie in $(14.13,14.14)$ and $(21,22)$, and the first four of $\beta$ in
$(6.02,6.03)$, $(10.24,10.25)$, $(12.98,13.00)$ and $(16.34,16.35)$, by the certified enclosures of FLINT
\cite{FLINTZetaZeros} and the LMFDB zero list of $\chi_{-4}$ with its reliability statement \cite{LMFDBDirichletBeta}. Both
spectra are symmetric and have no central zero. For $\gamma>0$ put $w(\gamma)=z_{1/2+i\gamma}^{20}=\exp(40i\arctan(1/(2\gamma)))$, and
take as targets $w_\zeta,w_\beta$ the values at the first positive zero of each function. For either target $w_*$,
\begin{equation}\label{rec:eq:boxes}
 \abs{1-w_*}^2>5/3,\qquad\abs{w-w_*}>3/25\quad\text{at every other folded zero of the two spectra}.
\end{equation}
In increasing order of positive ordinate the real parts of the first four folded frequencies of the union lie in
$(-.987,-.984)$, $(-.373,-.369)$, $(.030,.034)$ and $(.155,.157)$, the fourth belonging to $\zeta$, and all remaining
positive ordinates give real part greater than $.340$. Conjugate frequencies have the same real part, and the
absolute imaginary parts at the two targets exceed $.16$ and $.98$, which separates each from its conjugate. For
$\gamma\ge10$ the angle in the exponential decreases in $(0,\pi)$, so its cosine increases; on the first interval the angle
lies between $\pi$ and $3\pi/2$. Rational endpoint evaluations of $((\gamma^2-\frac14+i\gamma)/(\gamma^2+\frac14))^{20}$ prove all these
bounds.

With multiplicity in each sum, the two completed functional equations give
\begin{equation}\label{rec:eq:bothmass}
 \sum_{\rho\text{ of }\zeta}\frac1{\abs\rho^2}+\sum_{\rho\text{ of }\beta}\frac1{\abs\rho^2}=2-2\log\pi+2\frac{\beta'(1)}{\beta(1)}<\frac12 .
\end{equation}
Indeed, for either real completion $\Xi$ the difference of its logarithmic derivatives at $1$ and $0$ is $2\Xi'(1)/\Xi(1)$
by symmetry, and the Hadamard expansion gives the same difference as $\sum1/(\rho(1-\rho))=\sum1/\abs\rho^2$ by
\cref{src:thm:dgrh}. For $\zeta$ the expression is $2+\recgE-\log(4\pi)$ and for $\beta$ it is $\log(4/\pi)-\recgE+2\beta'(1)/\beta(1)$. For the
strict bound use $\beta(1)=\pi/4>0.78$, $\log\pi>1.14$ and $\beta'(1)<0.30$; the alternating series for $\beta'(1)$ is bounded above
by its five terms from $\log3/3$ through $\log11/11$, and termwise differentiation at $1$ follows from uniform
convergence of the alternating derivative series on a real neighbourhood of $1$.

Let $D(k)=A_{u,v}(k+1)-2A_{u,v}(k)+A_{u,v}(k-1)$ and $B=\max(\abs u,\abs v)>0$. The Fourier series of $D$ has an atom at
$w=z^{20}$ for each zero, with coefficient $-b_\rho z^r\abs{1-w}^2$, $b_\rho$ the scalar $u$ or $v$ times the multiplicity. By
$\abs{1-z^{20}}\le20\abs{1-z}$ and \eqref{rec:eq:bothmass} the absolute coefficients sum to less than $200B$, and the target
of a scalar of modulus $B$ has coefficient of modulus greater than $5B/3$ by \eqref{rec:eq:boxes}. Let
$K(X)=(64^{-1}\sum_{i=0}^{63}X^i)^4=\sum_{i=0}^{252}a_iX^i$, with nonnegative coefficients of sum one. In
$\sum_ia_i\bar w_*^{\,i}D(k_0+i)$ the target contributes its full coefficient times a unit factor, and every other frequency is
multiplied in modulus by at most $(25/96)^4$, by the geometric sum and \eqref{rec:eq:boxes}. Since $200(25/96)^4<1$, the
average has modulus greater than $2B/3$, so some $\abs{D(k_0+i)}>2B/3$. An affine function has zero second difference,
and the absolute coefficient sum of a second difference is $4$, so one of the $255$ indices from $k_0-1$ to $k_0+253$ has
$\abs{A_{u,v}-P}>B/6$. For two integer pairs the difference has $B\ge1$, and two fits with errors strictly below $1/12$
would contradict the inequality. The rational inequalities and frequency enclosures are checked in exact
arithmetic:
\begin{center}
\artifact{dirichlet_li_recovery/checks/window.txt}
\end{center}
\end{proof}

The count $255$ is sufficient; its optimality is not known.

\paragraph{Heights at sums of zero angles.}
Let $S_0=\set0$ and $S_j$ the set of sums modulo $2\pi$ of at most $j$ zero angles, a compact set. For
$y\in S_k\setminus S_{k-1}$ let $M_k(y)$ be the number of ordered $k$-tuples of zero occurrences with angle sum $y$,
a zero of multiplicity $m$ offering $m$ choices at each use.

\begin{theorem}[Multiplicities at every finite depth; \stG]\label{rec:thm:liecho}
For $y\in S_k\setminus S_{k-1}$, $k\ge1$, $M_k(y)$ is finite and positive, and uniformly for $\abs x\le x_0$,
\[
 \frac{p_{\chi,t}(y+ax)}{p_{\chi,t}(ax)}=\frac{M_k(y)}{k!}t^k+O_{\chi,y,x_0}(t^{k+1}),\qquad
 \frac{p_{\chi,t}(y+ax)}{t^kp_{\chi,t}(0)}\to\frac{M_k(y)}{k!(1+x^2)}.
\]
If $y\notin\bigcup_jS_j$ the ratio is $O_{\chi,y,x_0,N}(t^N)$ for every $N$. In particular, for every $y\ne0$,
\[
 \lim_{t\downarrow0}\frac{p_{\chi,t}(y)}{t\,p_{\chi,t}(0)}=\ord_{s=1/2+(i/2)\cot(y/2)}L(s,\chi),
\]
a central zero included at $y=\pi$.
\end{theorem}

\begin{proof}
Let $d=\operatorname{dist}(y,S_{k-1})>0$. A sum of $k$ angles within $d/2$ of $y$ has every summand of norm at least
$d/2$, so $M_k(y)$ is finite and $y$ is isolated in $S_k$; choose $h$ below a quarter of both distances. For
small $\eta>0$ let $\mu$ be the part of $\nu_\chi$ on angles of norm above $\eta$, $\ell=\mu(\mathbb T)$, and $g_t$ the density
of the remaining jumps; the exponents change by at most $2\ell$, so by \cref{rec:prop:liheat}
$\norm{g_t}_\infty\ll_\eta a^{-1}$ and $\norm{g_t-p_t}_\infty\ll_\eta t/a$. For the lifted small-jump process $Y_v$, with drift
$b_\eta$ and $V=\int\theta^2\dd\nu_\chi$, the exponential formula gives
$\mathbb E e^{\pm sY_v}\le\exp(vs\abs{b_\eta}+\tfrac12vs^2e^{s\eta}V)$; with $s=\log(1/t)/(2\eta)$ and $v=t/2$ this gives
$\Pr(\abs{Y_{t/2}}\ge h/2)\ll t^{h/(4\eta)}$, and splitting $g_t=g_{t/2}*g_{t/2}$ gives
$\sup_{\operatorname{dist}(z,0)\ge h}g_t(z)\ll_\eta a^{-1}t^{h/(4\eta)}$. Take $\eta<d/2$ with $h/(4\eta)>k+1$ and expand
$p_t=e^{-\ell t}\sum_{j\ge0}\frac{t^j}{j!}\mu^{*j}*g_t$. For $j<k$ and for nonmatching centres at $j=k$ the point $y+ax$ is at
distance at least $h$ from every centre; the matching centres contribute $M_k(y)t^kg_t(ax)/k!$; the terms
$j>k$ cost $O_\eta(t^{k+1}/a)$. Since $p_t(ax)\gg_{x_0}a^{-1}$, this proves the ratio, and \cref{rec:prop:liheat} the
profile. Off all $S_j$ use $\operatorname{dist}(y,S_N)>0$ in the same way.
\end{proof}

Hence $\lim\log(p_{\chi,t}(y)/p_{\chi,t}(0))/\log t$ is the least number of factors $1-1/\rho$ whose product is $e^{iy}$, or
$+\infty$. The points of finite depth form a countable dense set (multiples of a small angle form a fine
net), and every depth occurs, since equal consecutive $S_j$ would stabilize and a compact countable set
is not dense. The formula survives finite data: for fixed $k$ and $c>2k$, truncating the Fourier series
at $\abs n\le N$ with $t=c/N$ gives a trigonometric polynomial $P_N$ with $k!P_N(y)/(t^kP_N(0))\to M_k(y)$ for
$y\notin S_{k-1}$, $M_k(y)$ read as $0$ off $S_k$; with Hermitian coefficient errors $\eps_N$ rounding still succeeds
eventually if $\limsup\eps_N/t^{k-1}<1/(2\,k!)$, so the first layer tolerates any uniform coefficient error below
$\tfrac12$ (\stG). Indeed the omitted tail is $O(a^{-1}N^{-b'c})=o(t^k/a)$ for $b'<\tfrac12$ with $b'c>k$, and an
input error moves the polynomial by at most $(1+o(1))t\eps_Np_t(0)$.

\begin{theorem}[The typical height and the minimum; \stG]\label{rec:thm:litypical}
Let $L=\log(1/t)$. For small $t$, $\inf_{y\in\mathbb T}p_{\chi,t}(y)/p_{\chi,t}(0)\ge\exp(-L\log L/\log2-O_\chi(L))$, and for
Lebesgue-almost every $y\in\mathbb T$,
\[
 \lim_{t\downarrow0}\frac{-\log(p_{\chi,t}(y)/p_{\chi,t}(0))}{L\log L}=\frac1{\log2};
\]
the same limit holds for $\min_yp_{\chi,t}(y)$.
\end{theorem}

\begin{proof}
\emph{Lower bound.} The one-sided count gives a positive ordinate in every $[T,T+D_\chi]$ for large $T$, so,
as $\dd\theta/\dd T=-1/(T^2+\tfrac14)$, every $0<r\le r_0$ admits a positive angle in $[r-Cr^2,r]$, with $Cr_0<1$. After
a bounded number of copies of one fixed small angle, subtract at each step such an angle from the
residual; the residuals satisfy $Cr_j\le(Cr_0)^{2^j}$, so every $y$ is within $a$ of a sum of
$m\le\log L/\log2+O_\chi(1)$ angles, as $\log(1/a)=L+O(\log L)$. Delete one unit of intensity at each of the $d\le m$
distinct angles used; the remaining density $g_t$ satisfies $\norm{g_t-p_t}_\infty\ll a^{-1}tde^{2td}$, so $g_t\gg a^{-1}$ on
$\abs z\le a$, and the deleted atoms realize the required counts with probability at least $e^{-tm}t^m/m!$.
\emph{Upper bound.} The count gives $\nu_\chi\set{\abs u>r}\ll r^{-1}\log(e/r)$, and the $\delta$-neighbourhood of $S_j$ has
length $F_j(\delta)\le2(jr+\delta)+Cr^{-1}\log(e/r)F_{j-1}(\delta)$. With $\delta=e^{-\beta n}$, $k=\floor{(1-\eps)\log n/\log2}$ and
$r=\delta^{2^{-j}}$ at step $j$, $F_k(\delta)\le\exp(-\beta n^\eps/2)$, which is summable, so by Borel--Cantelli almost every $y$
has $\operatorname{dist}(y,S_k)>\delta$ for all large $n$. For $e^{-n-1}\le t\le e^{-n}$ truncate the jumps at $\eta=\delta/(16k)$; then
$t\nu_\chi\set{\abs u>\eta}\le\exp(-(1-\beta)n+O(\log n))$, the exponential moment bound with $s=(1-\beta)n/(2\eta)$ gives
$\sup_{\norm z\ge\delta}g_t(z)\ll a^{-1}e^{-s\delta/2}$ for the small-jump density, and
$ap_t(y)\ll e^{-4(1-\beta)nk}+(t\nu_\chi\set{\abs u>\eta})^{k+1}$. Dividing by $L\log L$ and letting $\beta,\eps\to0$ gives the
almost-everywhere limit; the minimum is squeezed between the lower bound and one typical point.
\end{proof}

Correspondingly, if $D_y(\delta)$ is the least number of zero angles whose sum is within $\delta$ of $y$, then
$\sup_yD_y(\delta)\le\log\log(1/\delta)/\log2+O_\chi(1)$ and $D_y(\delta)/\log\log(1/\delta)\to1/\log2$ for almost every $y$ and for
the supremum; the upper bound can be attained using only zeros with $\abs\gamma\le C_\chi\delta^{-1/2}$, finishing with one
positive and one negative angle of size about $\sqrt\delta$. Thus a point of the unit circle is approximated
within angular error $10^{-D}$ by a product of at most $\log_2D+O_\chi(1)$ factors $1-1/\rho$ of zeros of $L(s,\chi)$,
and for almost every target no fewer than $(1+o(1))\log_2D$ suffice; each step squares the residual.

\begin{example}[Thirteen asymmetric values for a quartic character]\label{rec:ex:liquartic}
Let $\chi$ be the quartic character modulo $5$ with $\chi(2)=i$. Interval arithmetic (Euler--Maclaurin with
explicit remainders, Rouch\'e's theorem on circles of radii $10^{-13}$ and $10^{-3}$, $301$ certified intervals
on the critical line, and a winding-number computation over $1918$ boundary segments of
$[-0.01,1.01]+i[-4.14,4.14]$) proves that $L(s,\bar\chi)$ has exactly one zero with $\abs{\Im s}\le4.14$ in the strip, simple,
on the critical line, with $4.13290370521275<\gamma_*<4.13290370521295$; its angle satisfies
$0.24079042878934<y_*<0.24079042878938$ and $13y_*<\pi<14y_*$. For $1\le k\le13$ a sum of at most $k$ angles of
$\bar\chi$ equals $ky_*$ only as $k$ copies of $y_*$, while the extremal angle of $\chi$ is $-y_*$, so
\[
 \frac{p_{\bar\chi,t}(ky_*)}{p_{\bar\chi,t}(0)}=\frac{t^k}{k!}+O_k(t^{k+1}),\qquad\frac{p_{\chi,t}(ky_*)}{p_{\chi,t}(0)}=O_k(t^{k+1}),
\]
with the roles exchanged at $-ky_*$; at $k=13$ the coefficient is $1/6227020800$. The strip count is
independent of \cref{src:thm:dgrh}:
\begin{center}
\artifact{dirichlet_li_heat_kernel/checks/expected.txt}\\
\artifact{dirichlet_li_heat_kernel/checks/rectangle_expected.txt}
\end{center}
\end{example}

\subsection{One path identifies the character}\label{li:sec:path}
\sources{grh_path_rigidity_2026-11-27}

Let $\chi$ be primitive of conductor $q_\chi$, with $q_\chi=1$ for $\zeta$, let $m_\chi(\gamma)$ be the
multiplicity of the ordinate $\gamma$ (zero if $\gamma$ is not an ordinate), and let $X^\chi$, $\nu_\chi$ and
$\Psi_\chi$ be the process, L\'evy measure and exponent of \cref{li:lem:levy} for the ordinates of
$L(s,\chi)$, counted with multiplicity and both signs. For $\gamma>0$ the jump rate at each of
$\pm\theta(\gamma)$ is $c_\chi(\gamma)=\frac12(m_\chi(\gamma)+m_\chi(-\gamma))$, and a central zero of multiplicity $m$ gives
rate $m/2$ at each of $\pm\pi$; for a real character the two multiplicities coincide and the rate is
$m_\chi(\gamma)$. Under \cref{src:thm:dgrh}, $\liEx e^{inX^\chi_t}=e^{-t\Re\lambda_\chi(n)}$ by \eqref{li:eq:liident}. Complex
conjugation reflects the ordinates and preserves $\nu_\chi$, so $X^\chi$ and $X^{\bar\chi}$ have the same law;
write $[\chi]=\set{\chi,\bar\chi}$.

For $0<\abs y<\pi$ put $G(y)=1/(2\tan(\abs y/2))$, so that $\pm\theta(\gamma)\mapsto\abs\gamma$, and for $T>1$ let
\[
 B_T=\set{y:T<G(y)\le2T}=\set{y:\theta(2T)\le\abs y<\theta(T)},
\]
a band of jump sizes asymptotic to $[1/(2T),1/T]$ with $\nu_\chi(B_T)<\infty$. The marked jumps
$\set{(s,\Delta X_s):0<s\le t,\ \Delta X_s\in B_T}$, $\Delta X_s=X_s-X_{s-}$, form a finite Poisson random measure with
intensity $\dd s\,\nu_\chi(\dd y)$. Hence for bounded real $f,g$ supported on $B_T$ the finite sums
$J_f=\sum_{0<s\le t}f(\Delta X_s)$ satisfy
\begin{equation}\label{li:eq:poisson}
 \liEx J_f=t\int f\,\dd\nu_\chi,\qquad\operatorname{Cov}(J_f,J_g)=t\int fg\,\dd\nu_\chi,
\end{equation}
and no compensation of smaller jumps is needed. Write $\Lambda(x)=\log p$ if $x=p^k$ is an integral prime
power and $\Lambda(x)=0$ otherwise, also for nonintegral $x$.

\begin{lemma}[Band masses and prime frequencies; \stG]\label{li:lem:bands}
As $T\to\infty$,
\[
 \nu_\chi(B_T)=\frac T\pi\log T+C_\chi T+O_\chi(\log T),\qquad C_\chi=\frac{\log(q_\chi/(2\pi))+2\log2-1}\pi.
\]
For fixed rational $x>1$,
\begin{equation}\label{li:eq:bandcos}
 \int_{B_T}\cos(G(y)\log x)\,\nu_\chi(\dd y)=-\frac{T\Lambda(x)\Re\chi(x)}{\pi\sqrt x}+O_{\chi,x}(\log^2T).
\end{equation}
Consequently, for fixed primes $p\ne r$,
\begin{gather*}
 V_{\chi,p}(T):=\int_{B_T}\cos^2(G(y)\log p)\,\nu_\chi(\dd y)=\frac T{2\pi}\log T+O_{\chi,p}(T),\\
 \int_{B_T}\cos(G(y)\log p)\cos(G(y)\log r)\,\nu_\chi(\dd y)=O_{\chi,p,r}(\log^2T).
\end{gather*}
\end{lemma}

\begin{proof}
The band mass is $N_\chi(2T)-N_\chi(T)$ with the count recalled after \cref{li:cor:universal}. The integral
in \eqref{li:eq:bandcos} is $\sum_{T<\abs\gamma\le2T}m_\chi(\gamma)\cos(\gamma\log x)$, the two half atoms of each ordinate
together carrying its multiplicity. Since $m_\chi(-\gamma)=m_{\bar\chi}(\gamma)$, it is the sum of the real parts of
$\sum_{T<\gamma\le2T}m_\xi(\gamma)x^{i\gamma}$ for $\xi=\chi,\bar\chi$, and \cref{zer:lem:rational} with $H=T$ gives
\eqref{li:eq:bandcos}; the estimate there covers $1<x<2$, which occurs below for ratios of primes. For
$V_{\chi,p}$ write $\cos^2v=\frac12(1+\cos2v)$ and use \eqref{li:eq:bandcos} at $x=p^2$. For the cross term the
product-to-sum identity leaves the frequencies $\log(pr)$ and $\abs{\log(p/r)}$; neither $pr$ nor a ratio of
distinct primes is a prime power, so both main terms vanish.
\end{proof}

For a path $X$ on $[0,t]$, a prime $p$ and $T>1$ define the jump statistic
\begin{equation}\label{li:eq:decoder}
 D_{T,t}(p;X)=-\frac{\pi\sqrt p}{tT\log p}\sum_{\substack{0<s\le t\\ \Delta X_s\in B_T}}\cos\bigl(G(\Delta X_s)\log p\bigr).
\end{equation}
It involves no conductor, character, zero location or multiplicity.

\begin{theorem}[Recovery of Euler coefficients from jumps; \stG]\label{li:thm:decoder}
For each primitive $\chi$, fixed $t>0$ and prime $p$,
\[
 \liEx_\chi D_{T,t}(p)=\Re\chi(p)+O_{\chi,p}\Bigl(\frac{\log^2T}T\Bigr),\qquad
 \operatorname{Var}_\chi D_{T,t}(p)=\frac{\pi p}{2t(\log p)^2}\,\frac{\log T}T+O_{\chi,p}\bigl((tT)^{-1}\bigr).
\]
Almost surely $D_{2^n,t}(p;X^\chi)\to\Re\chi(p)$ as $n\to\infty$, simultaneously for all primes. For fixed distinct
primes $p_1,\dots,p_k$,
\[
 \sqrt{\frac{tT}{\log T}}\bigl(D_{T,t}(p_j)-\Re\chi(p_j)\bigr)_{j=1}^k\Longrightarrow
 N_k\Bigl(0,\operatorname{diag}\Bigl(\frac{\pi p_j}{2(\log p_j)^2}\Bigr)_{j\le k}\Bigr),
\]
so the limiting errors at distinct primes are independent.
\end{theorem}

\begin{proof}
With $f=\cos(G\log p)\mathbf 1_{B_T}$, \eqref{li:eq:poisson} and \cref{li:lem:bands} give the mean and the variance.
Chebyshev's inequality gives
$\Pr(\abs{D_{2^n,t}(p)-\liEx D_{2^n,t}(p)}>\eps)\ll_{\chi,p,t,\eps}n2^{-n}$, which is summable; Borel--Cantelli and the mean give the almost sure limit, and one intersects over all primes and a
sequence of tolerances. For the joint law fix real $b_1,\dots,b_k$. The centred combination is a
compensated Poisson integral on $B_T$ with integrand
\[
 h_T=-(tT\log T)^{-1/2}\mathbf 1_{B_T}\sum_jb_j\frac{\pi\sqrt{p_j}}{\log p_j}\cos(G\log p_j),\qquad
 \norm{h_T}_\infty=O\bigl((tT\log T)^{-1/2}\bigr).
\]
Its log characteristic function at $u$ is
$t\int(e^{iuh_T}-1-iuh_T)\dd\nu_\chi=-\frac{u^2}2t\int h_T^2\dd\nu_\chi+R_T$ with
$\abs{R_T}\le\frac{\abs u^3}6\norm{h_T}_\infty t\int h_T^2\dd\nu_\chi$. By \cref{li:lem:bands},
$t\int h_T^2\dd\nu_\chi\to\sum_jb_j^2\pi p_j/(2(\log p_j)^2)$ and $R_T\to0$, while the bias times
$\sqrt{tT/\log T}$ tends to zero. The Cram\'er--Wold argument concludes.
\end{proof}

The independence is a limit theorem for Poisson jump counts; it neither uses nor implies linear
independence of ordinates, and multiplicities and shared ordinates are included in the rates. The
number of jumps in $B_T$ is of order $tT\log T$, the arithmetic signal of the unnormalized sum of order
$tT$, which explains the scale $\sqrt{\log T/(tT)}$.

\begin{proposition}[Finite-band confidence bound; \stU]\label{li:prop:confidence}
With $a_p=\pi\sqrt p/\log p$ and $V=V_{\chi,p}(T)$, for every $z>0$,
\[
 \Pr_\chi\bigl(\abs{D_{T,t}(p)-\liEx_\chi D_{T,t}(p)}\ge z\bigr)\le2\exp\Bigl(-\frac{tT^2z^2}{2a_p^2V+\frac23a_pTz}\Bigr).
\]
\end{proposition}

\begin{proof}
For $\abs f\le1$, comparing exponential series termwise gives
$\log\liEx e^{v(J_f-\liEx J_f)}=t\int(e^{vf}-1-vf)\dd\nu_\chi\le tv^2\int f^2\dd\nu_\chi/(2(1-\abs v/3))$ for $\abs v<3$. The
exponential Markov inequality, optimized in $v$, gives
$\Pr(\abs{J_f-\liEx J_f}\ge u)\le2\exp(-u^2/(2(tV+u/3)))$; substitute $u=tTz/a_p$.
\end{proof}

A confidence interval for $\Re\chi(p)$ must also contain the deterministic bias of
\cref{li:thm:decoder}.

\begin{lemma}[Real prime values determine the conjugacy orbit; \stU]\label{li:lem:orbit}
If primitive $\chi,\psi$ satisfy $\Re\chi(p)=\Re\psi(p)$ for all primes $p$, then $\psi\in[\chi]$; in particular
$q_\chi=q_\psi$.
\end{lemma}

\begin{proof}
Induce both to $(\Z/M\Z)^\times$, $M=\lcm(q_\chi,q_\psi)$. Every class contains a prime, so
$\chi+\bar\chi=\psi+\bar\psi$ there, and linear independence of characters gives equality of the multisets
$\set{\chi,\bar\chi}$ and $\set{\psi,\bar\psi}$ of induced characters. Primitive inducing characters are unique.
\end{proof}

Let $P_\chi^t$ be the law of $X^\chi$ on the Skorokhod space $D([0,t],\R)$ with its Borel $\sigma$-field.

\begin{theorem}[Path-law dichotomy and recovery; \stG]\label{li:thm:path}
For every $t>0$:
\begin{enumerate}[label=\textup{(\roman*)}]
\item $P_\chi^t=P_\psi^t$ if and only if $[\chi]=[\psi]$;
\item if $[\chi]\ne[\psi]$, then $P_\chi^t$ and $P_\psi^t$ are mutually singular;
\item one measurable map on $D([0,t],\R)$ returns $[\chi]$ almost surely under every $P_\chi^t$;
\item almost surely $[\chi]$ is recovered from every interval of positive length of one path, and from
the germ of the path at $0$.
\end{enumerate}
\end{theorem}

\begin{proof}
Conjugate characters have the same law. If the orbits differ, \cref{li:lem:orbit} gives a prime with
$\Re\chi(p)\ne\Re\psi(p)$, and the two limits of \cref{li:thm:decoder} define disjoint events of full measure
under the two laws. These events are measurable: on a band bounded away from $0$ a c\`adl\`ag path has
finitely many jumps, and integrals of Borel functions against its jump measure are measurable path
functionals. For (iii) enumerate the countably many orbits; the events on which every limit
exists and equals $\Re\chi(p)$ are pairwise disjoint by \cref{li:lem:orbit}; output the orbit on its event
and a default elsewhere. Stationary increments give the same decoder on every $[\alpha,\beta]$; intersect
over rational endpoints. For the germ let $F_n$ be the decoder on $[0,t/n]$ and $F$ its eventual value
(a default if none). For every $N$, $F$ is a function of $(F_n)_{n\ge N}$, hence measurable for the path on
$[0,t/N]$, and under each law all $F_n$ equal the orbit almost surely.
\end{proof}

Under each fixed law the recovered orbit is a constant, in agreement with the zero-one law; the theorem
compares different laws on one path space. Part (i) also follows from \cref{rec:thm:liheat}(ii), since
the law of $X^\chi_t$ modulo $2\pi$ has Fourier coefficients $e^{-t\Re\lambda_\chi(n)}$.

\begin{corollary}[Continuous perturbations and finitely many altered jumps; \stG]\label{li:cor:pathrobust}
The almost sure recovery of \cref{li:thm:path} persists after adding any continuous function to the
path and after altering, inserting or deleting finitely many jumps, even when these changes are chosen
after the path is observed.
\end{corollary}

\begin{proof}
A continuous perturbation changes no jump. Finitely many altered jumps have a positive least nonzero
size, so they leave $B_T$ for all large $T$, and the decoder sequences agree eventually.
\end{proof}

No stability under errors in every small jump is claimed: $G$ has derivative of order $T^2$ on $B_T$.

Fix nonconjugate primitive $\chi,\psi$ and a prime $p$ with $\Delta_p=\Re\chi(p)-\Re\psi(p)\ne0$. For $y$ in the
finite union of the atoms of $\nu_\chi,\nu_\psi$ in $B_T$ put $\lambda_y=\nu_\chi\set y$, $\mu_y=\nu_\psi\set y$, and
\[
 H_T(\chi,\psi)=\sum_{y\in B_T}\bigl(\sqrt{\lambda_y}-\sqrt{\mu_y}\bigr)^2,\qquad K_p=\frac{(\log p)^2\Delta_p^2}{2\pi p}.
\]

\begin{theorem}[Finite-band information and testing; \stG]\label{li:thm:information}
The Hellinger sum satisfies
\[
 \liminf_{T\to\infty}\frac{\log T}TH_T(\chi,\psi)\ge K_p.
\]
If the jumps in $B_T$ are observed during time $t$ with
their exact sizes, the least equal-prior error of a test between the two characters satisfies
\[
 P_{\rm err}(T,t)\le\tfrac12e^{-tH_T/2},\qquad
 P_{\rm err}(T,t)\le\tfrac12\exp\Bigl[-(1-o(1))\frac{t(\log p)^2\Delta_p^2}{4\pi p}\,\frac T{\log T}\Bigr].
\]
The counts at each size suffice; the jump times add no information.
\end{theorem}

\begin{proof}
With $f_y=\cos(G(y)\log p)$, \eqref{li:eq:bandcos} gives
$M_T:=\sum_yf_y(\lambda_y-\mu_y)=-\frac{T\log p}{\pi\sqrt p}\Delta_p+O(\log^2T)$. Cauchy--Schwarz and
$(\sqrt\lambda+\sqrt\mu)^2\le2(\lambda+\mu)$ give $M_T^2\le2H_T\int_{B_T}\cos^2(G\log p)\,\dd(\nu_\chi+\nu_\psi)$, and the last integral is
$\frac T\pi\log T+O(T)$ by \cref{li:lem:bands}. Under the two hypotheses the counts at the sizes $y$ are
independent Poisson variables with means $t\lambda_y$ and $t\mu_y$; the Hellinger affinity of Poisson laws with
means $a,b$ is $\exp(-(\sqrt a-\sqrt b)^2/2)$, so the product affinity is $e^{-tH_T/2}$, and the least error is half
the integral of the smaller density, at most half the affinity. Given the counts, the jump times are
uniform under both hypotheses.
\end{proof}

At a size threshold $\eps\asymp1/T$ the bound is of order $1/(\eps\log(1/\eps))$. It holds when the two
$L$-functions share many ordinates and is not claimed to be optimal, since atoms of only one character
add information.

\begin{theorem}[Finite-grid equivalence; \stU]\label{li:thm:grid}
For any primitive $\chi,\psi$ and fixed $0<t_1<\dots<t_k$, the laws of $(X^\chi_{t_1},\dots,X^\chi_{t_k})$ and
$(X^\psi_{t_1},\dots,X^\psi_{t_k})$ are mutually absolutely continuous.
\end{theorem}

\begin{proof}
By \cref{li:thm:allorders} each $X^\chi_s$, $s>0$, has a nonnegative entire density, not identically zero; its
real zeros are discrete, so it is positive almost everywhere and the law is equivalent to Lebesgue
measure. The independent increments have a product density positive almost everywhere, and the
triangular map from increments to positions preserves equivalence.
\end{proof}

\begin{corollary}[Grids against paths; \stG]\label{li:cor:grid}
For nonconjugate $\chi,\psi$ no classifier using a fixed finite time grid is almost surely correct under
both laws, whereas the complete path on the same interval admits one.
\end{corollary}

\begin{proof}
Such a classifier would give an event of probability one under one grid law and zero under the
other, contradicting \cref{li:thm:grid}; the second assertion is \cref{li:thm:path}.
\end{proof}

A c\`adl\`ag path carries infinitely many ever smaller jumps, and the null sets that separate the
path laws are not decided by finitely many coordinates. Likewise \cref{li:cor:universal} concerns
one-time densities and does not control the jump counts used by the decoder.

\begin{example}[Conductor five]\label{li:ex:five}
Let $\chi_2$ be the quadratic character modulo $5$ and $\chi_4$ the quartic one with $\chi_4(2)=i$; their orbits
differ. At $p=2$ the limits of \cref{li:thm:decoder} are $-1$ and $0$, the limiting variance is
$\pi/(\log2)^2$ under either law, $\liminf_T\frac{\log T}TH_T(\chi_2,\chi_4)\ge(\log2)^2/(4\pi)$, and
$P_{\rm err}(T,t)\le\frac12\exp[-(1-o(1))\frac{t(\log2)^2}{8\pi}\frac T{\log T}]$. Their scaled densities have the same
expansion $\sum_kQ_k/w_5(t)^k$ and differ by $O(\sqrt{a_5(t)})$ in every $C^r_b$ norm, their finite-grid laws are
equivalent, and their path laws are singular on every interval.
\end{example}

\begin{remark}[A finite Poisson model]\label{li:rem:illustration}
With the first $150$ positive zeros of $\zeta$, computed numerically and not certified, the finite
band $128<\gamma\le256$ ($T=128$) and $t=1$, the exact finite-model means of $D_{128,1}(p)$ are $0.97431$, $0.99960$
and $0.99115$ for $p=2,3,5$, with variances $0.15430$, $0.09322$ and $0.07055$; $50{,}000$ simulated
realizations agree with these moments. This illustrates the finite model and proves nothing about the
limits:
\begin{center}
\artifact{grh_path_rigidity_2026-11-27/checks/decoder_checks.json}
\end{center}
\end{remark}

\Cref{src:thm:dgrh} enters only through \eqref{li:eq:bandcos}: it turns $x^\rho$ in Landau's formula into the
phase $\sqrt x\,e^{i\gamma\log x}$, which is read from a jump size. The process, its densities and
\cref{li:thm:grid} use only the zero count. Banks studies differences of zero multisets of Dirichlet
$L$-functions through the same prime frequencies \cite{Banks2023Comparing}.

\subsection{The oriented kernel away from the origin}\label{li:sec:offorigin}
\sources{dirichlet-li-kernels}

Here $p_t=p_{\chi,t}$ is the oriented circle kernel of \cref{rec:prop:liheat}, with complex Li coefficients
$\lambda_\chi(n)$ and jump measure $\nu_\chi=\sum_\gamma\delta_{\theta_\gamma}$ on $\mathbb T=\R/2\pi\Z$, $L=\log(1/t)$ and $a=a(t)\sim\frac12tL$.
Let $S_m$ be the set of sums modulo $2\pi$ of at most $m$ zero angles, $M_k$ the ordered representation
counts, and $\delta_m(x)=\operatorname{dist}(x,S_m)$. At points of $S_k\setminus S_{k-1}$ the kernel is a rescaled copy of the
central Cauchy profile with weight $M_k/k!$, at points outside $\bigcup_mS_m$ it is $O(t^N)$ for every $N$
(\cref{rec:thm:liecho}), and at almost every point and at its minimum
$-\log p_t\sim L\log L/\log2$ (\cref{rec:thm:litypical}). The relative estimate of \cref{rec:thm:liecho} holds with
derivatives: for $x\in S_k\setminus S_{k-1}$ and fixed $R,r$,
\begin{equation}\label{li:eq:Crprofile}
 \Bigl\|ap_t(x+a\,\cdot)-\frac{M_k(x)}{k!}t^kap_t(a\,\cdot)\Bigr\|_{C^r([-R,R])}=O_{x,k,R,r}(t^{k+1})\qquad(\text{\stG}),
\end{equation}
and off $\bigcup_mS_m$ the left-hand rescaled density is $O(t^H)$ in $C^r([-R,R])$ for every $H$. Indeed
$a^{r+1}\sum_n\abs n^re^{-t\Re\lambda_\chi(n)}=O_r(1)$ by \eqref{rec:eq:lagarias}; removing the finitely many large jumps
multiplies the Fourier coefficients by factors of modulus at most $e^{2t\ell}$, so this bound and an $O(t)$
error in $C^r$ persist for the remaining density $g_t$; and the mass bound $O(t^{h/(4\eta)})$ for $g_t$ away from
$0$ becomes a $C^r$ bound after convolution with a mollifier of width $t^H$, which moves $r$th derivatives
by $O(t^H)$ and costs a factor $t^{-H(r+1)}$. The expansion of the proof of \cref{rec:thm:liecho} then
runs in $C^r$. Combined with \cref{li:prop:trace}, the profile at each minimal sum has the full
logarithmic expansion multiplied by $M_k(x)/k!$.

\paragraph{Exceptional decay rates.} Put
\[
 R_t(x)=\frac{-\log p_t(x)}{L\log L},\qquad F(x)=\liminf_{t\downarrow0}R_t(x),\qquad
 \omega(x)=\limsup_{m\to\infty}\frac{\log\log(1/\delta_m(x))}m,
\]
with $\omega(x)=\infty$ on $\bigcup_mS_m$. The greedy approximation after \cref{rec:thm:litypical} gives
$\delta_m(x)\le\exp(-c2^m)$ for large $m$, so $\omega\ge\log2$ everywhere. List the distinct values of $\abs{\theta_\gamma}$ in
nonincreasing order as $b_1\ge b_2\ge\cdots$. Every sum of zero angles is a signed sum
$\sum_{i\le k}e_ib_{j_i}$, $e_i=\pm1$.

\begin{lemma}[Separation from signed sums; \stU]\label{li:lem:separation}
Let $p,s>1$, $0<d\le1$, $pd>1$ and $0<\eta<1$ with $2\eta^d\zeta(pd)<1$. The set of $x\in\mathbb T$ lying in
infinitely many of the arcs of radius $\eta^k\prod_{i\le k}j_i^{-p}$ about the signed sums $\sum_{i\le k}e_ib_{j_i}$
($k\ge0$, all signs and indices) has zero $d$-dimensional Hausdorff measure. Every $x$ outside this set and
outside the countable set of signed sums satisfies, for all $m\ge1$,
\[
 \delta_m(x)\ge\exp\bigl\{-C_{x,s,p}(1+sp)^m\log(2m)\bigr\},\qquad\text{so}\quad\omega(x)\le\log(1+sp).
\]
\end{lemma}

\begin{proof}
The sum of the $d$th powers of the diameters of all arcs is at most $2^d\sum_k(2\eta^d\zeta(pd))^k<\infty$, so the
tails of the family are covers of vanishing mesh and vanishing $d$-cost. For $x$ outside the limsup set
and the centres, decreasing a constant handles the finitely many arcs containing $x$, and
$\operatorname{dist}(x,\sum e_ib_{j_i})\ge c_x\eta^k\prod j_i^{-p}$ for all tuples. The zero count gives
$\#\set{j:b_j\ge u}\ll u^{-1}\log(2/u)$, hence $b_j\le C_sj^{-1/s}$. Take a tuple of length $n\le m$ with magnitudes
in decreasing order, let $r_k$ be the distance from $x$ to the sum of its first $k$ signed terms,
$\delta=r_n$, and $B_k=\sum_{i\le k}\log j_i$. Stop at the first $k$ with $r_k\le2\delta$. Before it,
$r_k\le\delta+mb_{j_{k+1}}$, so $b_{j_{k+1}}>r_k/(2m)$, and the magnitude and separation bounds give
$B_{k+1}\le(1+sp)B_k+s(\log(2mC_s/c_x)+k\log(1/\eta))$. From $B_0=0$, geometric summation gives
$B_k\le C(1+sp)^k\log(2m)$ up to the stopping index, where $\delta\ge r_k/2\ge\frac12c_x\eta^k e^{-pB_k}$.
\end{proof}

With $d=1$ and $p,s$ near $1$ it gives $\omega\le\log2$ almost everywhere, the separation half of the almost-everywhere law of \cref{rec:thm:litypical}.

\begin{lemma}[Decay away from short sums; \stG]\label{li:lem:decay}
Let $c>0$ and $0<\beta<\beta'<1$, and for small $t$ put $m=\floor{c\log L}$, $\delta=\exp(-L^\beta)$ and $b=\exp(-L^{\beta'})$.
If $\delta_m(x)\ge\delta$, then
\[
 p_t(x)\le\frac Ca\Bigl[\exp\bigl(-\tfrac12e^{L^{\beta'}-L^\beta}+o(1)\bigr)+\exp\bigl(-(m+1)(L-O(L^{\beta'}))\bigr)\Bigr],
\]
with $C$ depending on $\chi,c,\beta,\beta'$ only; hence $R_t(x)\ge c+o(1)$ uniformly in such $x$.
\end{lemma}

\begin{proof}
Split the jumps at size $b$. The larger ones are compound Poisson of rate $\lambda\le V/b^2$,
$V=\int\theta^2\dd\nu_\chi$, so $t\lambda\le\exp(-L+O(L^{\beta'}))$. Let $q_v$ be the density of the remainder and $Y_v$ its
real lift; its Fourier coefficients are those of $p_v$ times factors of modulus at most $e^{2v\lambda}$, so
$\norm{q_v}_\infty=O(1/a)$ for $v=t,t/2$. With all jumps compensated the lift has drift $d_0$, and the
remainder has drift $d_b=d_0-\int_{\abs y>b}y\,\nu_\chi(\dd y)$, $\abs{d_b}\le\abs{d_0}+V/b$. Then
$\log\liEx e^{\pm Y_{t/2}/b}\le Ct(\abs{d_0}/b+V/b^2)=o(1)$. If $\operatorname{dist}(u,0)\ge\delta$, one of the two factors of
$q_t=q_{t/2}*q_{t/2}$ is evaluated at distance at least $\delta/2$ from $0$, and exponential Markov gives
$\sup_{\operatorname{dist}(u,0)\ge\delta}q_t(u)\le(C/a)\exp(-\delta/(2b)+o(1))$. On the event of at most $m$ large jumps their sum
lies in $S_m$, so this bound applies at the residual; otherwise bound the density by $C/a$ and use
$\Pr(N\ge m+1)\le(t\lambda)^{m+1}/(m+1)!$ for the Poisson number $N$ of large jumps. Since $\log(1/a)=O(L)$ the last
claim follows.
\end{proof}

\begin{theorem}[Spectrum of decay rates; \stG]\label{li:thm:spectrum}
\begin{enumerate}[label=\textup{(\roman*)}]
\item $F(x)=1/\omega(x)$ for every $x\in\mathbb T$, with $1/\infty=0$.
\item For $0<h\le1/\log2$, $\dim_{\mathrm H}\set{x:F(x)=h}=1/(e^{1/h}-1)$; the level $h=1/\log2$ has full measure.
\item $\set{x:F(x)=0}$ is a dense $G_\delta$ of Hausdorff dimension zero.
\item There is a dense $G_\delta$ set of points outside every $S_m$ at which $\liminf_{t\downarrow0}R_t(x)=0$ and
$\limsup_{t\downarrow0}R_t(x)=1/\log2$; at these points every value in $[0,1/\log2]$ is a subsequential limit of $R_t(x)$.
\end{enumerate}
\end{theorem}

\begin{proof}
(i) If $\omega(x)<B<\infty$, then $\delta_m(x)\ge\exp(-e^{Bm})$ for large $m$; for $c<1/B$ and $cB<\beta<\beta'<1$,
$m=\floor{c\log L}$ gives $e^{Bm}\le L^{cB}<L^\beta$, and \cref{li:lem:decay} gives $F(x)\ge c$; hence $F\ge1/\omega$. If
$B<\omega(x)$ there are arbitrarily large $m$ with $\delta_m(x)\le\exp(-e^{Bm})$; put $L=e^{Bm}$, $t=e^{-L}$, and take a sum
$z$ of at most $m$ angles with $\operatorname{dist}(x,z)\le t\le a$. Deleting one unit of intensity at each distinct angle
used, as in the proof of the lower bound of \cref{rec:thm:litypical}, gives $p_t(x)\ge c_0e^{-tm}t^m/(2a\,m!)$,
so $R_t(x)\le1/B+o(1)$ along these $t$ and $F(x)\le1/B$. Finally $F\ge0$ everywhere since $p_t\ll1/a$.

(ii) For $R>2$ let $E_R=\set{\omega\ge\log R}$ and $D=1/(R-1)$. For $D<d<1$ choose $p\in(1/d,R-1)$, $s>1$ with
$1+sp<R$, and $\eta$ with $2\eta^d\zeta(pd)<1$; by \cref{li:lem:separation}, $E_R$ has zero $d$-dimensional measure, so
$\dim_{\mathrm H}E_R\le D$. For the lower bound take small $r_0>0$ and $r_{m+1}=r_m^R$. By the one-sided zero
count every small $r$ admits a positive angle in $[r-Cr^2,r]$, so one can choose $M_m\asymp1/r_m$ angles in
$[r_m/4,3r_m/4]$ with mutual distances at least $cr_m^2$. Starting from $[0,r_0]$, replace each interval
$[z,z+r_m]$ of level $m$, $z$ a sum of $m$ angles, by the children $[z+y,z+y+r_{m+1}]$ for the chosen $y$; they
are disjoint since $R>2$. Every limit point $x$ has $\delta_m(x)\le r_m$, so lies in $E_R$. Give the children
equal shares of mass and let $A_m=\prod_{i<m}M_i$; then $\log A_m=D\log(1/r_m)+O(m+1)$. An interval of length
$r_{m+1}\le\ell\le r_m$ meets at most two intervals of level $m$ and at most $C(1+\ell/r_m^2)$ children of each, so has
mass at most $CA_m^{-1}\min(1,r_m+\ell/r_m)$, which is $O(\ell^{d'})$ for every $d'<D$ because $RD=D+1$ and
$A_m^{-1}=O(r_m^{d'})$. By Frostman's lemma the measure vanishes on sets of dimension below $D$, in particular on
$\bigcup_jE_{R+1/j}$, so $\set{\omega=\log R}$ has dimension $D$. Put $R=e^{1/h}$ and use (i); for $h=1/\log2$ use
\cref{rec:thm:litypical}.

(iii) The set $\set{\omega=\infty}$ lies in every $E_R$, so has dimension zero, and equals
$\bigcap_{n,M}\bigcup_{m\ge M}\set{x:\delta_m(x)<\exp(-e^{nm})}$; each inner union is open and contains the dense set
$\bigcup_mS_m$, and Baire's theorem applies.

(iv) At a point $z$ of some $S_k$, $p_t(z)\gg_zt^k/a$, so $R_t(z)\to0$; points with $R_t\to1/\log2$ have full
measure, hence are dense. For $j\ge1$ and $N\ge3$ the sets $\bigcup_{n\ge N}\set{R_{e^{-n}}<1/j}$ and
$\bigcup_{n\ge N}\set{R_{e^{-n}}>1/\log2-1/j}$ are open, by continuity and positivity of $p_t$, and dense. Intersect them
and remove the countable set $\bigcup_mS_m$. The bounds $p_t\ll1/a$ and the uniform lower bound of
\cref{rec:thm:litypical} give the two extreme limits over all real $t$, and $R_t(x)$ is continuous in $t$, so
every intermediate value is a subsequential limit.
\end{proof}

For example the rate $1/\log3$ occurs on a set of dimension $\frac12$. Picard studied small-time densities
of L\'evy processes at points reachable by finitely many jumps and at their limits; his dyadic model
has a different almost-everywhere logarithmic rate \cite{Picard1997Density,Picard1997Accessible}.

\begin{example}[Two neighbouring two-jump peaks for $\zeta$; \stG]\label{li:ex:zetapeaks}
Let $a_j=2\arctan(1/(2\gamma_j))$ for the $j$th positive ordinate of $\zeta$; the first eight zeros are simple.
Since $44a_1<\pi$, for $1\le k\le44$ the two largest points of $S_k\setminus S_{k-1}$ in $(0,\pi)$ are $ka_1$ and
$(k-1)a_1+a_2$, with $M_k=1$ and $M_k=k$, so the coefficients $M_k/k!$ are $1/k!$ and $1/(k-1)!$. Rational
enclosures from Odlyzko's ordinates give $0.00003542639447445315<v-u<0.00003542639447445316$ for $u=2a_2$ and
$v=a_1+a_7$, both above $a_1$,
with $a_1+a_8<u<v$; the largest two-jump sum without $a_1$ is $2a_2$. Hence $u$ and $v$ are consecutive
two-jump centres with $M_2(u)=1$, $M_2(v)=2$, and the peak at $v$ is asymptotically twice as high as the
peak at $u$:
\begin{center}
\artifact{dirichlet-li-kernels/zeta-check.txt}
\end{center}
\end{example}

\subsection{What one time slice remembers}\label{li:sec:slice}
\sources{dirichlet-li-kernels; dirichlet_li_recovery}

Write $\mathbb T$ multiplicatively as $\set{z\in\C:\abs z=1}$ in this subsection. For $z=e^{i\theta}\in\mathbb T\setminus\set1$
the point $\rho=1/(1-z)=\frac12+i\gamma$ lies on the critical line, with \emph{ordinate}
$\gamma=\gamma(z)=\frac12\cot(\theta/2)$; this identifies the critical line with the punctured circle, and both the
multiplicity and the sign of $\gamma$ matter. An \emph{admissible multiset} $E$ in $\mathbb T\setminus\set1$ has
nonnegative integral multiplicities $m_E(z)$ with $\sum_zm_E(z)\abs{1-z}^2<\infty$ and with
$\lim_{\eps\downarrow0}\sum_{\abs{1-z}\ge\eps}m_E(z)\Im z$ existing. Put
\[
 \lambda_E(n)=\lim_{\eps\downarrow0}\sum_{\abs{1-z}\ge\eps}m_E(z)(1-z^n),\qquad H_{E,t}(n)=e^{-t\lambda_E(n)}\qquad(n\in\Z,\ t>0);
\]
the limit exists since $1-z^n+in\Im z=O_n(\abs{1-z}^2)$, and $H_{E,t}$ is the Fourier sequence of the weak limit of
the truncated compound Poisson laws, which fixes the drift. The \emph{time slice} of $E$ at time $t$ is the
sequence $H_{E,t}$, equivalently that probability measure on $\mathbb T$; the positive modes determine the
negative ones by conjugation, and $H_{E,t}(0)=1$. A \emph{collision} at time $t$ is a pair of distinct admissible
multisets with the same time slice. For a finitely supported integral measure $b=\sum_zb_z\delta_z$ write
$\widehat b(n)=\sum_zb_zz^n$ and $\norm b_1=\sum_z\abs{b_z}$; its positive and negative parts are admissible. Put
$\beta_t=2\pi i/t$, $\zeta_k=e^{2\pi i/k}$, $\mu_k=\set{z:z^k=1}$, $U_k=(\Z/k\Z)^\times$ and $\varphi(k)=\abs{U_k}$, and
for $a\in U_k$ let $\sigma_a$ be the automorphism $\zeta_k\mapsto\zeta_k^a$ of $\Q(\zeta_k)$; $\mu(\cdot)$ is the M\"obius
function. Everything down to \cref{li:rem:fibersizes} is a statement about arbitrary admissible multisets and
uses no information on zeros.

\begin{lemma}[Almost periodic second differences; \stU]\label{li:lem:per}
For admissible $E,F$ put $D=\lambda_E-\lambda_F$ and $u=\Delta^2D$, $\Delta^2f(n)=f(n+1)-2f(n)+f(n-1)$. Then
$u(n)=\int z^n\dd\tau(z)$ with $\tau=\sum_z\abs{1-z}^2(m_E(z)-m_F(z))\delta_z$ of finite variation, and for each $\delta>0$ the
$\delta$-almost periods of $u$ have bounded gaps. If $e^{-tu(n)}=1$ for all $n$, $\tau$ is supported on a finite cyclic
subgroup of $\mathbb T$. If $e^{-tu(n)}\ne1$ for some $n$, there are $\eta>0$ and $L_0$ such that every block of $L_0$
consecutive integers contains $n$ with $\abs{e^{-tu(n)}-1}>\eta$.
\end{lemma}

\begin{proof}
The formula for $u$ holds at each cutoff and passes to the limit. Approximate $u$ within $\delta/4$ by
$\sum_jc_jz_j^n$; in the closure of the orbit of $(z_1,\dots,z_J)$ in the torus take an identity neighbourhood on
which $\sum_j\abs{c_j}\abs{z_j^h-1}<\delta/2$; finitely many orbit translates cover the closure, which gives returns at
bounded distance from every integer, and each return is a $\delta$-almost period. If $e^{-tu}\equiv1$, $u$ takes values
in $\beta_t\Z$, so a positive $\abs{\beta_t}/2$-almost period $r$ is an exact period, all Fourier coefficients
of $(z^r-1)\tau$ vanish and $\tau$ lives on the $r$th roots of unity. The last claim follows from uniform continuity
of the exponential on the bounded range of $u$ and the bounded gaps of almost periods.
\end{proof}

\begin{lemma}[Finite torsion reduction; \stU]\label{li:lem:torsion}
Let $E\ne F$ be admissible with $H_{E,t}=H_{F,t}$, and $b=\sum_z(m_E(z)-m_F(z))\delta_z$. Then $b$ is supported on
$\mu_q$ for some $q$, has mass $\sum_zb_z=0$, and $\lambda_E-\lambda_F=-\widehat b$ with $\widehat b(n)\in\beta_t\Z$ for every
$n\in\Z$; moreover $\beta_t^2\in\Q$. Conversely, if an integral measure $b$ on $\mu_q\setminus\set1$ has mass zero and
$\widehat b(n)\in\beta_t\Z$ for all $n\ge1$, then adding its positive and its negative part to a common admissible
multiset gives two multisets with the same time slice.
\end{lemma}

\begin{proof}
Equality of the slices gives $D(n)=\lambda_E(n)-\lambda_F(n)\in\beta_t\Z$, hence $e^{-t\Delta^2D}\equiv1$, and
\cref{li:lem:per}, with $\abs{1-z}>0$ for $z\ne1$, confines $b$ to some $\mu_q$. Writing $s=\sum_zb_z$ gives
$D(n)=s-\widehat b(n)$. Its real part vanishes since $\beta_t$ is imaginary, so Fourier inversion on $\mu_q$ gives
$b_z+b_{\bar z}=2s\mathbf 1_{z=1}$; as $b_1=0$, $s=0$ and $D=-\widehat b$. Since $E\ne F$, some $D(n_0)$ is a nonzero
multiple of $\beta_t$, hence $\beta_t\in\Q(\zeta_q)$. For $a\in U_q$ the automorphism $\sigma_a$ sends $D(n_0)$ to $D(an_0)$,
another nonzero integral multiple of $\beta_t$. Thus $a\mapsto\sigma_a(\beta_t)/\beta_t\in\Q^\times$ is a character of a finite
group, with values $\pm1$, and every automorphism fixes $\beta_t^2$. Conversely $D=-\widehat b$ takes values in $\beta_t\Z$,
the negative modes by conjugation, so $e^{-tD}\equiv1$.
\end{proof}

Call $t$ \emph{exceptional} if $\beta_t^2\in\Q$, that is, $(\pi/t)^2\in\Q$. Then $\beta_t^2<0$; let $d<0$ be the
fundamental discriminant of $\Q(\beta_t)$, $f=\abs d$, and $\chi_d=(d/\cdot)$ the primitive odd quadratic character of
conductor $f$, extended by zero to nonunits. Its Gauss sum satisfies
\[
 G=\sum_{j\bmod f}\chi_d(j)\zeta_f^j,\qquad G^2=d,\qquad\sum_{j\bmod f}\chi_d(j)\zeta_f^{jn}=G\chi_d(n)\quad(n\in\Z),
\]
including the zero at nonunits \cite{Davenport2000Multiplicative}, and $\abs G=\sqrt f$. Since $\beta_t/G\in\Q^\times$,
define coprime positive integers $A,B$ by $\abs{\beta_t/G}=A/B$. For $m\ge1$ put
$\Omega_m=\sum_{j\bmod fm}\chi_d(j)\delta_{\zeta_{fm}^j}$.

\begin{theorem}[Determining times and all invisible changes; \stU]\label{li:thm:times}
Fix $t>0$. If $(\pi/t)^2\notin\Q$, the map $E\mapsto H_{E,t}$ on admissible multisets is injective. If $t$ is
exceptional, the differences $m_E-m_F$ of the pairs with $H_{E,t}=H_{F,t}$ form exactly the group
\begin{equation}\label{li:eq:lattice}
 \liKk_t=\bigoplus_{m\ge1}\Z\,\frac A{\gcd(A,m)}\,\Omega_m ,
\end{equation}
and any common admissible multiset may be added to both members of a pair. The elements of $\liKk_t$ supported
on $\mu_q$ are spanned by the generators with $fm\mid q$; their rank is zero if $f\nmid q$ and equals the number of
divisors of $q/f$ otherwise. In particular two multisets with the same slice agree, with multiplicity,
outside the roots of unity of some finite order, and every exceptional time has collisions.
\end{theorem}

\begin{proof}
Off the exceptional times \cref{li:lem:torsion} gives injectivity. Let $t$ be exceptional and $b$ a collision
difference on $\mu_q$, and take $Q=\lcm(q,f)$. For $a\in U_Q$, $\sigma_a(G)=\chi_d(a)G$, and $\beta_t/G$ is rational, so
$\widehat b(an)=\sigma_a(\widehat b(n))=\chi_d(a)\widehat b(n)$ by \cref{li:lem:torsion}, and Fourier inversion gives
$b_{z^a}=\chi_d(a)b_z$. On the shell of primitive $k$th roots, $k\mid Q$, the power action of $U_Q$ is transitive. A
nonzero coefficient there forces $\chi_d$ to be trivial on the kernel of $U_Q\to U_k$, equivalently $f\mid k$ by the
definition of the conductor; the coefficients on such a shell have the form $v_k\chi_d(j)$ at $\zeta_k^j$, $j\in U_k$, with
$v_k\in\Z$. Each $\Omega_m$ is equivariant under this character; on its highest shell $k=fm$ its coefficient at $\zeta_{fm}$
is $1$, and its other shells have orders properly dividing $fm$. Removing shell coefficients in decreasing order
gives a unique integral expansion $b=\sum_{fm\mid q}c_m\Omega_m$, and the argument applies to every integral
equivariant array. Writing $j=u+f\ell$ and summing first over $\ell$ gives
\begin{equation}\label{li:eq:DFT}
 \widehat\Omega_m(n)=\begin{cases}mG\chi_d(n/m),&m\mid n,\\0,&m\nmid n.\end{cases}
\end{equation}
Each $\Omega_m$ has mass zero, so by \cref{li:lem:torsion} equality of the slices is equivalent to $\widehat b(n)\in\beta_t\Z$
for $n\ge1$. The integer $\widehat b(n)/G=\sum_{m\mid n}mc_m\chi_d(n/m)$ lies in $(A/B)\Z$ exactly when $A$ divides it. At
$n=m$ the new term is $mc_m$, so induction gives $A\mid mc_m$ for all $m$, and the same formula proves
sufficiency. Finally $A\mid mc_m$ exactly when $A/\gcd(A,m)\mid c_m$.
\end{proof}

The denominator $B$ drops out: exceptional times with the same field and numerator $A$ have identical
invisible changes. At $t=\pi$ one has $f=4$, $G=2i$, $A=1$ and $\Omega_1=\delta_i-\delta_{-i}$, so single atoms at $i$ and at $-i$
give the same law; the examples of Hasebe and Hotta are of this kind \cite{HasebeHotta2022Additive}. At $t=\pi/3$ in
the same field $A=3$, the first two generators are $3\Omega_1$ and $3\Omega_2$, and $\Omega_3$ itself is invisible. Exact
quadratic Gauss sums for $d=-3,-4,-7,-8$ are checked in \artifact{dirichlet-li-kernels/gauss-check.txt}, and the
classification is compared with exact cyclotomic Fourier computations on bounded integral arrays for
$q\in\set{3,4,5,6,8,12}$ in
\artifact{dirichlet_li_recovery/checks/heat.txt}. The splitting of primitive shells into the two classes of a
quadratic character also occurs in the classification of integral mixed circulant graphs
\cite{KadyanBhattacharjya2023Integral}; here every Fourier coefficient must in addition lie in the prescribed lattice
$\beta_t\Z$, whose spacing fixes the least mass of an invisible change.

\paragraph{Low-ordinate probes and the sharp band.}
For $k\ge3$ let $\Gamma_k=\set{\frac12\cot(\pi j/k):j\in U_k}$ be the ordinates of the primitive $k$th roots of unity
and $h_k=\min_{\gamma\in\Gamma_k}\abs\gamma$; these are the locations available to a competing multiset.

\begin{theorem}[Exact supplementary data; \stU]\label{li:thm:probe}
At an exceptional time, the time slice together with the multiplicities at all ordinates in a set $W\subseteq\R$
determines every admissible multiset exactly when
\begin{equation}\label{li:eq:hit}
 W\cap\Gamma_k\ne\emptyset\qquad\text{for every positive multiple $k$ of $f$}.
\end{equation}
For $W=[-r,r]$ the exact threshold, endpoint included, is $r_f=\frac12\tan(\pi/f)$. The positive ordinates $h_k$, one
for each multiple $k$ of $f$, suffice; they accumulate only at $0$.
\end{theorem}

\begin{proof}
A collision difference has shell coefficients $v_k\chi_d(j)$, as in the proof of \cref{li:thm:times}, so if it vanishes at
one primitive $k$th root it vanishes on the whole shell. Conversely, if $W$ misses the shell of order $k$, assign
$\chi_d(j)$ on that shell and zero elsewhere; this integral equivariant array is an integral combination of the
$\Omega_m$, so $A$ times it is invisible by \eqref{li:eq:lattice}, and its positive and negative parts give a collision
agreeing on $W$. A unit $j$ nearest to $k/2$ minimizes the absolute height, which gives
\begin{equation}\label{li:eq:heights}
 h_k=\begin{cases}\frac12\tan(\pi/(2k)),&k\text{ odd},\\ \frac12\tan(\pi/k),&4\mid k,\\ \frac12\tan(2\pi/k),&k\equiv2\pmod4,\ k\ge6:\end{cases}
\end{equation}
for odd $k$ choose $j=(k-1)/2$; if $4\mid k$, $k/2-1$ is odd and its common divisors with $k$ divide $2$; for
$k\equiv2\pmod4$ the neighbours of $k/2$ are even, while $k/2-2$ is odd with common divisors dividing $4$. If $f$ is
even then $4\mid f$, and the largest $h_k$ occurs at $k=f$; if $f$ is odd it occurs at $k=2f$, since odd multiples
have tangent argument at most $\pi/(2f)$, multiples divisible by four at most $\pi/(4f)$, and the remaining multiples
at most $\pi/f$. Both cases give the maximum $r_f$, and \eqref{li:eq:heights} gives $h_k\to0$.
\end{proof}

\begin{corollary}[Universal central band; \stU]\label{li:cor:band}
At any $t>0$, two admissible multisets with equal time slices and equal multiplicities at every ordinate
$\abs\gamma\le\sqrt3/2$ coincide. The constant is sharp over all times: at $t=2\pi/\sqrt3$ the single atoms at
$e^{\pm i\pi/3}$, the zeros $\frac12\pm i\sqrt3/2$, have the same time slice.
\end{corollary}

\begin{proof}
Off the exceptional times use \cref{li:thm:times}. At an exceptional time $f\ge3$, and $r_f$ decreases in $f$, so
$r_f\le r_3=\sqrt3/2$ and \cref{li:thm:probe} applies. At $t=2\pi/\sqrt3$, $\beta_t=i\sqrt3$, and the difference
$\delta_{e^{i\pi/3}}-\delta_{e^{-i\pi/3}}$ has transform $2i\sin(n\pi/3)\in i\sqrt3\,\Z$, so \cref{li:lem:torsion} applies.
\end{proof}

\begin{corollary}[Ambiguity below the threshold; \stU]\label{li:cor:rank}
At an exceptional time and for $r>0$, the invisible changes that vanish on all ordinates in $[-r,r]$ form a group
of finite rank
\[
 R_f(r)=\#\set{k\in f\Z_{>0}:h_k>r}=\frac\pi{2fr}+O(1)\qquad(r\downarrow0).
\]
\end{corollary}

\begin{proof}
Each missed shell allows one integral amplitude, and $A$ times each individual shell is in the kernel. Put
$u=\arctan(2r)$. For even $f$, $k=fm$ is missed precisely when $m<\pi/(fu)$. For odd $f$ the tests are $m<\pi/(2fu)$
for odd $m$, $m<\pi/(fu)$ for $4\mid m$, and $m<2\pi/(fu)$ for $m\equiv2\pmod4$, on residue classes of densities
$\frac12,\frac14,\frac14$. Either parity gives $\pi/(fu)+O(1)$, and $u=2r+O(r^3)$.
\end{proof}

An admissible multiset has finitely many atoms in a bounded interval of ordinates, because $\abs{1-z}$ is bounded
away from zero there. The data in the band can therefore be supplied as a finite list of signed ordinates with
multiplicities, together with the information that the list is complete.

\paragraph{Least number of altered atoms.}
$\norm b_1$ counts the atoms left on the two sides of a collision after common multiplicities are cancelled; by
\cref{li:lem:torsion} each side carries $\norm b_1/2$ of them.

\begin{lemma}[Norm of a character projection; \stU]\label{li:lem:norm}
Let $f\mid Q$ and let $w\in\mu_Q$ have order $k$. The average $P(w)=\varphi(Q)^{-1}\sum_{a\in U_Q}\chi_d(a)w^a$ vanishes unless $f\mid k$.
For $k=fh$ and $w=\zeta_k^v$, $v\in U_k$,
\begin{equation}\label{li:eq:projection}
 P(w)=\frac{\chi_d(v)\mu(h)\chi_d(h)G}{\varphi(k)},\qquad\abs{P(w)}\le\frac{\sqrt f}{\varphi(f)},
\end{equation}
with equality in the bound precisely at order $k=f$, or at $k=2f$ when $f$ is odd.
\end{lemma}

\begin{proof}
If $\chi_d$ is nontrivial on the kernel of $U_Q\to U_k$, the sum over that kernel vanishes. Otherwise $f\mid k$, equal
fibres reduce the average to $U_k$, and changing variables by $v$ extracts $\chi_d(v)$. Insert
$\mathbf 1_{(a,k)=1}=\sum_{e\mid(a,k)}\mu(e)$. Terms with $(e,f)>1$ vanish; for $e\mid k$ coprime to $f$ the inner sum is
$\chi_d(e)\sum_{b\bmod k/e}\chi_d(b)\exp(2\pi ib/(k/e))$, so $e\mid h$, and summing over lifts modulo $f$ gives zero unless $k/e=f$,
by a geometric series. The surviving term is $\mu(h)\chi_d(h)G$, read as zero when $(h,f)>1$. Since $\varphi(k)\ge\varphi(f)$
the bound follows; equality requires $(h,f)=1$ and $\varphi(k)=\varphi(f)$, hence $\varphi(h)=1$, and precisely $h=1$ and, for odd
$f$, $h=2$ qualify.
\end{proof}

\begin{theorem}[Exact distance; \stU]\label{li:thm:dist}
At an exceptional time the least value of $\norm b_1$ over nonzero $b\in\liKk_t$ is $A\varphi(f)$. Both $A\Omega_1$ and $\Omega_A$
attain it, and $\Omega_A$ has no repeated atoms on either side. For packets of at most $M$ atoms added to an
arbitrary fixed background, the time slice is injective exactly when $M<A\varphi(f)/2$.
\end{theorem}

\begin{proof}
Choose $n$ with $\widehat b(n)\ne0$. By \eqref{li:eq:DFT} and \cref{li:thm:times}, $\widehat b(n)/G$ is a nonzero multiple of $A$, so
$\abs{\widehat b(n)}\ge A\sqrt f$. Since $\widehat b(an)=\chi_d(a)\widehat b(n)$, averaging over $a\in U_Q$ gives
$\abs{\widehat b(n)}=\abs{\sum_zb_zP(z^n)}\le\norm b_1\sqrt f/\varphi(f)$ by \cref{li:lem:norm}. Each $\Omega_m$ has $m\varphi(f)$ coefficients of
modulus one, so the examples attain equality. Differences of packets of at most $M$ atoms have norm at most
$2M$; at the threshold take the positive and negative parts of a minimizer, since a common background multiplies
both slices by the same nonzero sequence.
\end{proof}

Let $q$ be the least positive integer with $\supp b\subseteq\mu_q$, so that the support generates $\mu_q$, and write
$q=fm$.

\begin{theorem}[All collisions at the minimum distance; \stU]\label{li:thm:eq}
Let $t$ be exceptional and $b\in\liKk_t$ with $\norm b_1=A\varphi(f)$. If $f$ is even or $m$ is odd, then $m\mid A$ and
$b=\pm(A/m)\Omega_m$. If $f$ is odd and $m$ is even, put $s=\chi_d(2)$; then $b=c\,\Omega_m+e\,\Omega_{m/2}$ with
$c=(Y-sX)/m\in\Z\setminus\set0$ and $e=2X/m\in\Z$ for some $X,Y\in\set{-A,0,A}$, not both zero, and every such choice
gives a minimizer with least order $q$. All minimizers are supported on $\mu_{fA}$ for even $f$ and on $\mu_{2fA}$ for
odd $f$.
\end{theorem}

\begin{proof}
Equality in both bounds of the preceding proof forces $\abs{\widehat b(n)}=A\sqrt f$ for every $n$ with $\widehat b(n)\ne0$, and
$\abs{P(z^n)}=\sqrt f/\varphi(f)$ at every $z\in\supp b$. Thus each $z^n$ has order $f$, or possibly $2f$ for odd $f$. The support
generates $\mu_q$, so its $n$th powers generate a group of order $q/\gcd(q,n)$, which is $f$ or $2f$; on the frequency
side only the shell of order $f$ and, when $f$ is odd and $2f\mid q$, the shell of order $2f$ occur. On each frequency
shell $\widehat b/G$ transforms by $\chi_d$ under multiplication of the index by units modulo $q$, so it is determined by its
value at one representative. The representative of order $f$ is $n=m$, where $\widehat\Omega_m/G$ equals $m$ and vanishes off that
shell by \eqref{li:eq:DFT}. When this is the only shell, inversion gives $b=\pm(A/m)\Omega_m$ and integrality forces $m\mid A$.

In the second case put $X=\widehat b(m/2)/G$ and $Y=\widehat b(m)/G$, the values at the representatives of orders $2f$ and
$f$. By \eqref{li:eq:DFT}, $\widehat\Omega_m/G$ takes the values $0,m$ there and $\widehat\Omega_{m/2}/G$ the values $m/2,sm/2$, and both vanish
off these shells. Solving at the two representatives and inverting gives the stated $c,e$, integral by the
triangular basis of \cref{li:thm:times}, and $c\ne0$ by minimality of $q$. Conversely these conditions give the
Fourier lattice, as $X,Y$ are multiples of $A$. At $\zeta_q^j$ the coefficient of $b$ is $(sc+e)\chi_d(\ell)$ for even $j=2\ell$
and $c\chi_d(j)$ for odd $j$, and each parity has $(m/2)\varphi(f)$ positions coprime to $f$, so
\[
 \norm b_1=\tfrac12m\varphi(f)(\abs c+\abs{sc+e})=\tfrac12\varphi(f)(\abs{Y-sX}+\abs{Y+sX})=\varphi(f)\max(\abs X,\abs Y).
\]
The coefficient at $\zeta_q$ is $c\ne0$, so the least order is $q$. Finally $mc=Y-sX$ is a nonzero element of
$\set{\pm A,\pm2A}$, so $m\mid2A$.
\end{proof}

\begin{example}[An odd conductor]\label{li:ex:oddcond}
At $t=2\pi/\sqrt3$ one has $f=3$, $A=1$ and $s=-1$. At $q=3$ the minimal pair consists of the single atoms at
$\zeta_3^{\pm1}$. At $q=6$ take $X=Y=1$, so $c=e=1$: the measure $\Omega_2+\Omega_1$ cancels on the primitive third roots and
has coefficients $+1,-1$ at $\zeta_6,\bar\zeta_6$, which is the pair of \cref{li:cor:band}. Up to exchanging the sides and
adding a common background, \cref{li:thm:eq} gives no further minimal pairs at this time.
\end{example}

\paragraph{Positivity and the number of possible spectra.}

\begin{theorem}[Every time slice has a finite fibre; \stU]\label{li:thm:fiber}
At an exceptional time, for each multiple $k$ of $f$ let $S_k$ be the total multiplicity of $E$ at the ordinates
$+h_k$ and $-h_k$. These integers are determined by the moduli $\abs{H_{E,t}(n)}$, only finitely many are nonzero,
and
\[
 \#\set{F\text{ admissible}:H_{F,t}=H_{E,t}}\le\prod_{f\mid k}(1+S_k)<\infty .
\]
When $A=1$ equality holds for every prescribed finite list of positive $S_k$, all other $S_k$ being zero. In
particular a multiset with no atom of ordinate $\abs\gamma\le r_f$ is determined by its time slice among all
admissible multisets, and the radius $r_f$ is sharp for this guarantee.
\end{theorem}

\begin{proof}
The moduli give $\Re\lambda_E(n)=-t^{-1}\log\abs{H_{E,t}(n)}$, whose central second differences are the Fourier coefficients
of the finite measure $\frac12\sum_z\abs{1-z}^2(m_E(z)+m_E(\bar z))\delta_z$; Fourier uniqueness recovers every total on a
conjugate pair. The pairs at $\pm h_k$ are distinct for different $k$, since they lie on distinct primitive shells, and
they lie in $[-r_f,r_f]$, where an admissible multiset has finitely many atoms. The multiplicity at $+h_k$ has at most
$1+S_k$ possible values, and once these are chosen \cref{li:thm:probe} determines the whole multiset; when all $S_k=0$
this is uniqueness without supplementary data. For equality when $A=1$, on each prescribed shell put
multiplicity $S_k$ at every point with $\chi_d(j)=+1$ and zero where $\chi_d(j)=-1$; the two members of a conjugate pair have
opposite character signs. The array $\chi_d(j)$ confined to one shell is invisible by \cref{li:thm:times}, and
independently transferring $0,1,\dots,S_k$ units from the positive half to the negative half gives the stated number of
distinct multisets with one slice. For any $A$, the single-shell construction in the proof of \cref{li:thm:probe} at a
shell with $h_k=r_f$ gives two multisets with one slice, both avoiding every smaller closed band.
\end{proof}

\begin{remark}[Every fibre size occurs]\label{li:rem:fibersizes}
At every exceptional time every positive integer is the size of some fibre. Take $m=A^2$ and $k=fm$. The array $V_k$
equal to $\chi_d(j)$ on the primitive $k$th roots and zero elsewhere satisfies, by inclusion--exclusion on the exponents
divisible by primes of $m$ outside $f$,
\[
 V_k=\sum_{e\mid m,\ (e,f)=1}\mu(e)\chi_d(e)\,\Omega_{m/e}.
\]
For each squarefree $e\mid A^2$ one has $A\mid m/e$, so every term lies in $\liKk_t$, and the single-shell construction of
the preceding proof gives exactly $S_k+1$ multisets for any $S_k\ge0$.
\end{remark}

\paragraph{Arithmetic multisets.}
For a finite product $P$ of primitive Dirichlet $L$-functions, the empty product and $\zeta$ included, let
$E_P$ be the union of the Cayley multisets $\set{1-1/\rho}$ of its factors. By \cref{src:thm:dgrh}, $E_P$ lies in
$\mathbb T\setminus\set1$ (with $-1$ for a central zero), the cutoff $\abs{1-z}\ge\eps$ means $\abs\rho\le1/\eps$, $\sum\abs\rho^{-2}<\infty$, and
starred convergence of the Li zero sums \cite{Lagarias2007Li} makes $E_P$ admissible with $\lambda_{E_P}$ the complex Li
coefficients of \cref{rec:sec:liheat}. Write $H_{P,t}=H_{E_P,t}$.

\begin{corollary}[Products of $\zeta$ and $\beta$ against arbitrary multisets; \stG]\label{li:cor:zetabetaslice}
Let $\beta(s)=L(s,\chi_{-4})$. For every $t>0$ and every finite product $P=\zeta^{a_1}\beta^{a_2}$, $E_P$ is the only admissible
multiset with time slice $H_{P,t}$, even when the competitor is not assumed to come from an $L$-function.
\end{corollary}

\begin{proof}
Both zero sets are symmetric, with no central zero, and their least positive ordinates lie in $(14.13,14.14)$
and $(6.02,6.03)$ (see \cref{rec:thm:li255}); so $E_P$ has no atom of ordinate $\abs\gamma\le6$, and $6>\sqrt3/2\ge r_f$. Off the exceptional times use \cref{li:thm:times}, and at
the exceptional ones \cref{li:thm:fiber}.
\end{proof}

\begin{corollary}[Every time determines a primitive product; \stG]\label{li:thm:products}
Let $t>0$ and let $P,Q$ be finite products of primitive Dirichlet $L$-functions. If $H_{P,t}(n)/H_{Q,t}(n)$
tends to a finite nonzero limit as $n\to+\infty$, then $P$ and $Q$ have the same primitive factors with
multiplicity, and the limit is $1$. More precisely, with $C_{P,t}(n)=H_{P,t}(n+1)H_{P,t}(n-1)/H_{P,t}(n)^2$, if the factor
multisets differ there are $\eta>0$ and $L_0\ge1$ such that every block of $L_0$ consecutive positive integers
contains $n$ with $\abs{C_{P,t}(n)/C_{Q,t}(n)-1}>\eta$.
\end{corollary}

\begin{proof}
Let $\mathcal C$ be the set of characters occurring in $P$ or $Q$, with multiplicities $a_\chi$ and $a'_\chi$, and $b=t(a-a')$.
Then $H_{P,t}/H_{Q,t}=\exp(-A_b)$ in the notation of \cref{rec:thm:liblock} with $h=1$ and $r=0$. If the ratio tends to
$\ell\ne0$, the constant calibration $\log\ell$ makes the left side of \eqref{rec:eq:calibration} tend to zero along
intervals of fixed length moving to infinity, so $b=0$ by \cref{rec:thm:liblock} with $d=0$ and $M=\norm b_2$; then the
ratio is identically $1$. If the factors differ, $b\ne0$, and $C_{P,t}(n)/C_{Q,t}(n)=\exp(-\Delta^2A_b(n-1))$ with the
forward difference of \cref{rec:thm:liblock}(ii), which gives the block statement.
\end{proof}

Thus no two distinct primitive products collide at any time, although every exceptional time has
collisions among admissible multisets (\cref{li:thm:times}). The multiset of $\bar P$ is the conjugate of $E_P$, hence
$\abs{H_{P,t}}^2=H_{P\bar P,t}$. So a finite nonzero limit of $\abs{H_{P,t}(n)}/\abs{H_{Q,t}(n)}$ determines exactly the total
multiplicity of the factors in each conjugacy orbit $\set{\chi,\bar\chi}$, and the symmetric law is blind precisely to
conjugating individual factors. For one character this is \cref{rec:thm:liheat}, which also allows different
times, a rotation and agreement only in natural density; finite blocks on an arithmetic progression of indices
suffice by \cref{rec:thm:liblock}.

\subsection{A small-deviation law for sums over zeros}\label{ana:sec:smalldev}
{\raggedright\sources{rh_sampling_and_small_deviations_2026-11-26; dirichlet_spectral_conditioning_2026-11-27; spectral_extremes_2026-11-27}}

Sums over zeros with weights $1/\gamma^2$ also occur as random variables. Let $\gamma_j>0$ be a
multiset with $N(T)=\#\set{j:\gamma_j\le T}\sim T\log T/(2\pi)$, and suppose that for
$c,D,q>0$ and $\beta\in\R$
\[
 H(u)=c\prod_j\Bigl(1+\frac{u^2}{\gamma_j^2}\Bigr),\qquad
 \log\frac{H(u)}c=\frac u2\Bigl(\log\frac{qu}{2\pi}-1\Bigr)+\beta\log u+\log\frac Dc-\frac1{48u}+R_H(u)
\]
as $u\to+\infty$, with $R_H^{(k)}(u)=O(u^{-2-k})$ for $k=0,1,2$. For $\theta>0$ let $G_j$ be
independent gamma variables of shape $\theta$ and mean $\theta$, and $X_\theta=\sum_jG_j/\gamma_j^2$.

\begin{theorem}[Small deviations of gamma sums; \stU]\label{ana:thm:gammasum}
For small $x>0$ let $v=-W_{-1}(-8\pi x/(\theta q))$ and $u=\theta v/(4x)$, so that
$v=\log(qu/2\pi)>1$, and put
\[
 A_\theta(x)=\frac{2(c/D)^\theta}{\sqrt{\pi\theta}}\frac{u^{-\theta\beta-1/2}}{\sqrt{v-1}}
 \exp\Bigl(-\frac{\theta u(v-2)}4\Bigr).
\]
Then $X_\theta$ has a smooth density $f_\theta$, and as $x\downarrow0$
\[
 \anaPr(X_\theta\le x)=A_\theta(x)\Bigl(1+\frac\theta{48u}+O\Bigl(\frac1{u\log u}\Bigr)\Bigr),\qquad
 f_\theta(x)=u^2A_\theta(x)\Bigl(1+\frac\theta{48u}+O\Bigl(\frac1{u\log u}\Bigr)\Bigr).
\]
In particular $x\log\anaPr(X_\theta\le x)/\log^2(1/x)\to-\theta^2/16$.
\end{theorem}

\begin{proof}
\emph{Tilting.} With $K(t)=\theta\sum_j\log(1+t/\gamma_j^2)$ one has
$\anaEx e^{-tX_\theta}=e^{-K(t)}=(c/H(\sqrt t))^\theta$. Under
$\dd\anaPr_t=e^{-tX_\theta+K(t)}\dd\anaPr$ the summands are independent gamma variables with
rates $\gamma_j^2+t$, of total mean $m(t)=K'(t)$ and variance $\sigma^2(t)=-K''(t)$, and the
hypothesis gives, with $t=u^2$ and $\ell=\log(qu/2\pi)$,
\[
 m(u^2)=\theta\Bigl(\frac\ell{4u}+\frac\beta{2u^2}+\frac1{96u^3}+O(u^{-4})\Bigr),\qquad
 \sigma^2(u^2)=\theta\Bigl(\frac{\ell-1}{8u^3}+\frac\beta{2u^4}+O(u^{-5})\Bigr).
\]

\emph{Local limit.} Let $p_t$ be the $\anaPr_t$-density of $(X_\theta-m(t))/\sigma(t)$,
$n=u\log u$, $b_j=[\sigma(t)(\gamma_j^2+t)]^{-1}$ and
$\kappa_3(t)=2\theta\sum_jb_j^3=O(n^{-1/2})$. Then $\theta\sum b_j^2=1$ and $\max b_j=O(n^{-1/2})$;
the at least $a_1n$ ordinates in $[u,2u]$ have $b_j^2\ge a_2/n$, so the characteristic function
satisfies $\abs{\phi_t(w)}\le(1+a_3w^2/n)^{-a_4n}$. For $\abs w\le n^{1/12}$,
$\log\phi_t(w)=-w^2/2-i\kappa_3w^3/6+O(w^4/n)$, and the tails are $O(e^{-a_7n^{1/6}})$. Fourier
inversion gives $\sup_y\abs{p_t(y)-\varphi(y)-\frac{\kappa_3}6(y^3-3y)\varphi(y)}=O(n^{-1})$.

\emph{Saddle point.} Let $m(t_*)=x$, $u_*=\sqrt{t_*}$ and $a=t_*\sigma(t_*)\asymp\sqrt n$. The
exact identities
\[
 \anaPr(X_\theta\le x)=\frac{e^{t_*x-K(t_*)}}a\int_0^\infty e^{-z}p_{t_*}(-z/a)\,\dd z,\qquad
 f_\theta(x)=\frac{e^{t_*x-K(t_*)}p_{t_*}(0)}{\sigma(t_*)}
\]
give both quantities as $e^{t_*x-K(t_*)}/\sqrt{2\pi\sigma^2(t_*)}$, times $1/t_*$ for the
distribution function, with relative error $O(1/(u_*\log u_*))$.

\emph{Evaluation.} The explicit $u,v$ solve $x=\theta v/(4u)$, and the expansion of $m$ gives
$u_*-u=O(1/\log u)$. For $\Phi(w)=xw^2-\theta\log(H(w)/c)$ one has
$\Phi'(u)=-\theta\beta/u+O(u^{-2})$ and $\Phi''=O(\log u/u)$ near $u$, so
$\Phi(u_*)-\Phi(u)=O(1/(u\log u))$, and
$u^2\sqrt{2\pi\sigma^2(u^2)}=\frac{\sqrt{\pi\theta}}2u^{1/2}\sqrt{v-1}(1+O(1/(u\log u)))$.
Finally the hypothesis gives
$e^{\Phi(u)}=(c/D)^\theta u^{-\theta\beta}\exp(-\theta u(v-2)/4)(1+\theta/(48u)+O(u^{-2}))$.
Smoothness of the density follows from the polynomial decay of the characteristic function.
\end{proof}

\begin{theorem}[Small deviations from zeta and real characters; \stG]\label{ana:thm:smalldev}
Let $\gamma_j$ be the positive ordinates of the nontrivial zeros of $\zeta$, with multiplicity,
$E_j$ independent exponential variables of mean $1$, $X=\sum_jE_j/\gamma_j^2$ and
$c_0=\xi(1/2)>0$. With $v=-W_{-1}(-8\pi x)$ and $u=v/(4x)$,
\[
 \anaPr(X\le x)=\frac{2^{5/4}c_0}{\pi^{3/4}}\frac{u^{-9/4}}{\sqrt{v-1}}e^{-u(v-2)/4}
 \Bigl(1+\frac1{48u}+O\Bigl(\frac1{u\log u}\Bigr)\Bigr)\qquad(x\downarrow0),
\]
the density is $u^2$ times the same expression, and
$x\log\anaPr(X\le x)/\log^2(1/x)\to-1/16$. More generally \cref{ana:thm:gammasum} holds for
$\zeta$ with $q=1$, $\beta=7/4$, $D=(\pi/2)^{1/4}$, $c=c_0$, and for a primitive real
nonprincipal character $\chi$ of conductor $q$, with $\chi(-1)=(-1)^a$, central order $m$ and
the ordinates of its noncentral zeros, with
\[
 c=c_\chi=\frac{\Lambda_\chi^{(m)}(1/2)}{m!}>0,\qquad\beta=\frac{2a-1}4-m,\qquad
 D=2\sqrt\pi\Bigl(\frac q{2\pi}\Bigr)^{(2a+1)/4},
\]
where $\Lambda_\chi(s)=(q/\pi)^{(s+a)/2}\Gamma((s+a)/2)L(s,\chi)$.
\end{theorem}

\begin{proof}
By \cref{src:thm:dgrh} the paired Hadamard product gives
$\xi(\frac12+u)=c_0\prod_j(1+u^2/\gamma_j^2)$, the Riemann--von Mangoldt formula gives the
count, and Stirling's expansion with $\log\zeta(u+\frac12)=O(2^{-u})$ gives
$\log\xi(\frac12+u)=\frac u2(\log\frac u{2\pi}-1)+\frac74\log u+\frac14\log\frac\pi2-\frac1{48u}+O(u^{-2})$,
with differentiable remainder. For $\chi$, the functional equation has a real root number
$\eps$ with $\eps(-1)^m=1$, so $H_\chi(u)=\Lambda_\chi(\frac12+u)/u^m$ is even and entire of order
one, equal to $c_\chi\prod_j(1+u^2/\gamma_j^2)$. Since $L(s,\chi)>0$ for real $s>1$ and has no
real zero in $s>1/2$, $c_\chi>0$. With $b=(2a+1)/4$, Stirling's expansion gives
$\log H_\chi(u)=\frac u2(\log\frac{qu}{2\pi}-1)+(b-\frac12-m)\log u+\log(2\sqrt\pi)+b\log\frac q{2\pi}
+(b^2-b+\frac16)/u+O(u^{-2})$, and $b^2-b+\frac16=-\frac1{48}$ for $b=\frac14,\frac34$.
\end{proof}

The reciprocal-$\xi$ law itself is classical \cite{RoynetteYor2005Wald}, and Hadamard
factorization methods for such laws are in \cite{GaoHannigLeeTorcaso2003Laplace}; a tilted
local Gaussian limit for the same density appears in \cite{Sra2026Positivity}. No simplicity
of zeros, pair correlation or central nonvanishing is used.

\paragraph{Conditioning a spectral gamma sum to be small.}
\Cref{ana:thm:gammasum} gives the size of the lower tail. Its conditional structure is as
follows: the small value is shared among the spectral coordinates along a deterministic profile,
with Brownian-bridge fluctuations, an independent exponential slack at the boundary, and a sharp
threshold for the number of low coordinates that look like independent gamma variables, while the
largest individual terms form a marked Poisson process at a lower height, independent of the whole
bulk path. The
probability statements hold for every regularly varying spectrum; \cref{src:thm:dgrh} enters only
to identify the law with a reciprocal completed $L$-function.

Let $0<\gamma_1\le\gamma_2\le\cdots\to\infty$, with repetitions, and $N(T)=\#\set{j:\gamma_j\le T}$.
Assume that $N$ is regularly varying of index $\alpha>0$, that is $N(rT)/N(T)\to r^\alpha$ for every
$r>0$. Fix $p>\alpha$ and $\theta>0$, let $G_j$ be independent gamma variables of shape $\theta$ and
rate $1$, and put
\[
 X=\sum_j\frac{G_j}{\gamma_j^p},\qquad K(t)=\theta\sum_j\log\Bigl(1+\frac t{\gamma_j^p}\Bigr),
\]
so that $\anaEx e^{-tX}=e^{-K(t)}$. For $0<x<\anaEx X$ let $t=t_x$ be the unique positive root of
$K'(t)=\theta\sum_j(\gamma_j^p+t)^{-1}=x$, and $u=t^{1/p}$; this \emph{exact} saddle replaces the
explicit $u$ of \cref{ana:thm:gammasum}, and $t,u\to\infty$ as $x\downarrow0$. For $c\in[0,\infty]$
put
\[
 X_x(c)=\sum_{\gamma_j\le cu}\frac{G_j}{\gamma_j^p},\quad
 m_x(c)=\theta\sum_{\gamma_j\le cu}\frac1{\gamma_j^p+t},\quad
 \sigma_x^2=\theta\sum_j\frac1{(\gamma_j^p+t)^2},\quad Y_x=t(x-X).
\]
Let $C_q=\int_0^\infty\alpha r^{\alpha-1}(1+r^p)^{-q}\,\dd r$ for $q>\alpha/p$, so that
$C_1=\alpha\pi/(p\sin(\pi\alpha/p))$ and $C_2=(1-\alpha/p)C_1$, and
\[
 A(c)=\frac1{C_1}\int_0^c\frac{\alpha r^{\alpha-1}}{1+r^p}\,\dd r,\qquad
 V(c)=\frac1{C_2}\int_0^c\frac{\alpha r^{\alpha-1}}{(1+r^p)^2}\,\dd r .
\]
Let $B$ be a standard Brownian bridge, $\Phi$ the standard normal distribution function and
$\varphi_\lambda$ the centred normal density of variance $\lambda$. Total variation is
$d_{\mathrm{TV}}(P,Q)=\sup_E\abs{P(E)-Q(E)}$.

\begin{theorem}[Joint conditional spectral law; \stU]\label{ana:thm:condlaw}
\begin{enumerate}[label=\textup{(\roman*)}]
\item $tx\sim\theta C_1N(u)$, $t^2\sigma_x^2\sim\theta C_2N(u)$ and
 $\sigma_x/x\sim\sqrt{C_2}/(\sqrt\theta\,C_1\sqrt{N(u)})$.
\item Fix $0<c_1<\dots<c_k<\infty$ and $m\ge0$. Conditionally on $X\le x$ the vector
\[
 \Bigl(\Bigl(\frac{X_x(c_i)-m_x(c_i)}{\sigma_x}\Bigr)_{i\le k},\ Y_x,\
 \Bigl(\frac{tG_j}{\gamma_j^p}\Bigr)_{j\le m}\Bigr)
\]
converges in total variation to $\bigl((B(V(c_i)))_{i\le k},E,(\Gamma_j)_{j\le m}\bigr)$, where $E$ is
exponential of rate $1$, the $\Gamma_j$ are gamma of shape $\theta$ and rate $1$, and the three
objects are independent.
\item Conditionally on $X\le x$, $\sup_{0\le c\le\infty}\abs{X_x(c)/x-A(c)}\to0$ in probability, and
the same holds with $X$ in place of $x$.
\end{enumerate}
\end{theorem}

The centring $m_x(c)$ is the exact tilted mean; replacing it by its leading term is in general
false on the fluctuation scale (\cref{ana:rem:centring}). At $c=\infty$ the standardized value is
$-Y_x/(t\sigma_x)$, so total variation is claimed only at finite cutoffs. Convergence in law of the
whole bulk path, jointly with the largest individual terms, is \cref{ana:thm:functional}; it is not a
total-variation statement, since the finite paths jump and the limit is continuous.

\begin{lemma}[Spectral sums; \stU]\label{ana:lem:specsums}
For $q>\alpha/p$, $N(u)^{-1}\sum_j(1+(\gamma_j/u)^p)^{-q}\to C_q$, also with sum and integral restricted
to a fixed interval $(a,b]$, $0\le a<b\le\infty$. Moreover $\sum_j\gamma_j^{-p}<\infty$ and
$\sum_{\gamma_j\le h}\gamma_j^{2p}\sim\frac\alpha{\alpha+2p}h^{2p}N(h)$ as $h\to\infty$.
\end{lemma}

\begin{proof}
The measures $N(u)^{-1}\sum_j\delta_{\gamma_j/u}$ converge to $\alpha r^{\alpha-1}\dd r$ on compact
subintervals of $(0,\infty)$; near $0$ the integrand is at most $1$ and the mass below $\delta$ tends
to $\delta^\alpha$. At infinity, given $\eps>0$ one has $N(2T)/N(T)\le2^{\alpha+\eps}$ for $T\ge T_0$,
hence $N(ru)/N(u)\le C_\eps r^{\alpha+\eps}$ for $u\ge T_0$, $r\ge1$; with $\eps<pq-\alpha$ the part with
$\gamma_j/u>R$ is at most $C\sum_{\ell\ge0}(2^\ell R)^{-pq}(2^{\ell+1}R)^{\alpha+\eps}$, uniformly small.
The same dyadic estimate gives $\sum_j\gamma_j^{-p}<\infty$. Finally
$\sum_{\gamma_j\le h}\gamma_j^{2p}=h^{2p}N(h)-2p\int_0^hr^{2p-1}N(r)\,\dd r$, and dominated convergence on
$[0,1]$ after the substitution $r=hs$ gives the limit $1-2p/(\alpha+2p)$. The values of $C_1,C_2$ are
beta integrals.
\end{proof}

Under the tilted measure $\dd\anaPr_t=e^{-tX+K(t)}\dd\anaPr$ the variables
$\widetilde G_j=(1+t/\gamma_j^p)G_j$ are independent gamma of shape $\theta$ and rate $1$ (first for
finite groups, by integrating out the independent tail). Put
\[
 a_j=\frac1{1+(\gamma_j/u)^p},\qquad\varsigma^2=\theta\sum_ja_j^2=t^2\sigma_x^2,\qquad
 S=tX=\sum_ja_j\widetilde G_j,
\]
so $S$ has $\anaPr_t$-mean $tx$ and variance $\varsigma^2$. \Cref{ana:lem:specsums} gives (i) of
\cref{ana:thm:condlaw}. The event $X\le x$ has positive probability: take a finite initial segment
whose expected tail is below $x/4$, apply Markov's inequality to the tail and make the finitely many
initial summands small.

\begin{lemma}[Uniform local limit on a spectral band; \stU]\label{ana:lem:band}
Let $I=(a,b]$ with $0\le a<b\le\infty$, delete any fixed finite set of indices, and put
$Z_{I,u}=\varsigma^{-1}\sum_{\gamma_j/u\in I}a_j(\widetilde G_j-\theta)$. For large $u$ it has a
continuous density $f_{I,u}$, and $f_{I,u}\to\varphi_{\lambda_I}$ uniformly and in $L^1$, where
$\lambda_I=C_2^{-1}\int_I\alpha r^{\alpha-1}(1+r^p)^{-2}\,\dd r$. If a single index $j$ is deleted, the
convergence is uniform also in $j$. The same holds for the whole spectrum
with a lower block $\set{\gamma_j\le h(u)}$ removed, if $h=o(u)$, with limiting variance $1$.
\end{lemma}

\begin{proof}
By \cref{ana:lem:specsums}, $\theta\sum_Ia_j^2/\varsigma^2\to\lambda_I$, and since $0<a_j\le1$,
\[
 \log\phi_{I,u}(z)=\theta\sum_I\Bigl[-\frac{iza_j}\varsigma-\log\Bigl(1-\frac{iza_j}\varsigma\Bigr)\Bigr]
 =-\frac{\theta z^2}{2\varsigma^2}\sum_Ia_j^2+O_{\theta,z}(\varsigma^{-1})\to-\frac{\lambda_Iz^2}2 .
\]
For a compact $[a',b']$ inside $I$ with $a'>0$ there are at least $dN(u)$ indices there, each with
$a_j\ge(1+b'^p)^{-1}$; as $\varsigma^2\asymp N(u)=n$,
$\abs{\phi_{I,u}(z)}=\prod_I(1+z^2a_j^2/\varsigma^2)^{-\theta/2}\le(1+Az^2/n)^{-Dn}$, unchanged by a finite
deletion. This is at most $\exp(-ADz^2/(1+A))$ for $\abs z\le\sqrt n$, and
$\int_{\sqrt n}^\infty(1+Az^2/n)^{-Dn}\dd z\le C\sqrt n(1+A)^{-Dn/2}\to0$. Hence the characteristic
functions converge in $L^1$, Fourier inversion gives uniform convergence of densities, and Scheff\'e's
lemma gives $L^1$ convergence. Deleting one index changes the variance fraction by at most
$\theta/\varsigma^2$ and leaves at least $dN(u)-1$ indices in $[a',b']$, whatever the index, so all
these bounds are uniform in it. For a deleted block, $h\le\delta u$ eventually gives
$\limsup N(h)/N(u)\le\delta^\alpha$, so the removed variance $\theta N(h)$ is $o(\varsigma^2)$, and the
interval $[u,2u]$ still supplies the envelope.
\end{proof}

\begin{proof}[Proof of \cref{ana:thm:condlaw}]
Under $\anaPr_t$ let $W=(S-tx)/\varsigma$ and $D_u=\anaEx_t[e^{\varsigma W}\anaone_{W\le0}]$. By
\cref{ana:lem:band} the density $f_u$ of $W$ tends uniformly to $\varphi_1$, so
\begin{equation}\label{ana:eq:Du}
 \varsigma D_u=\int_0^\infty e^{-y}f_u(-y/\varsigma)\,\dd y\to\varphi_1(0)=(2\pi)^{-1/2},
\end{equation}
and the density of $\anaPr(\cdot\mid X\le x)$ with respect to $\anaPr_t$ is exactly
$D_u^{-1}e^{\varsigma W}\anaone_{W\le0}$. Omit the $m$ low coordinates and split the rest into the bands
$I_1=(0,c_1],\dots,I_k=(c_{k-1},c_k],I_{k+1}=(c_k,\infty)$, with independent standardized sums of
densities $f_{i,u}$ and limiting variances $\lambda_i=V(c_i)-V(c_{i-1})$ ($c_0=0$), $\lambda_{k+1}=1-V(c_k)$.
Let $U_i$ be the standardized increment on the full band $I_i$. For values $r=(r_1,\dots,r_m)>0$ of
$tG_j/\gamma_j^p=a_j\widetilde G_j$ put $\delta_u(r)=\varsigma^{-1}\sum_{j\le m}(r_j-\theta a_j)$; their
tilted density $q_u(r)=\prod_j r_j^{\theta-1}e^{-r_j/a_j}/(\Gamma(\theta)a_j^\theta)$ tends to the
product $q(r)$ of unit-rate gamma densities, and $\delta_u(r)\to0$. Changing the last increment to
$y=-\varsigma\sum_{i\le k+1}U_i$ (Jacobian $1/\varsigma$), the conditional density of
$(U_1,\dots,U_k,Y_x,r)$ is
\[
 \frac{e^{-y}\anaone_{y\ge0}q_u(r)}{\varsigma D_u}\,f_{1,u}(z_1-\delta_u(r))\prod_{i=2}^kf_{i,u}(z_i)\,
 f_{k+1,u}\Bigl(-\frac y\varsigma-\sum_{i\le k}z_i\Bigr),
\]
which converges almost everywhere to
$e^{-y}\anaone_{y\ge0}q(r)\varphi_1(0)^{-1}\prod_{i\le k}\varphi_{\lambda_i}(z_i)\,
\varphi_{\lambda_{k+1}}(-\sum_iz_i)$. The normal part divided by $\varphi_1(0)$ is the density of
independent centred normal increments of variances $\lambda_i$ conditioned to sum to zero, so the limit
integrates to one and Scheff\'e's lemma gives total-variation convergence. The cumulative increments
have covariance $V(\min(c_a,c_b))-V(c_a)V(c_b)$, which is the law of $(B(V(c_i)))_i$; an invertible
linear map preserves total variation, and the factorization gives independence.

For (iii), $m_x(c)/x\to A(c)$ by \cref{ana:lem:specsums} and $\sigma_x/x\to0$, so (ii) gives
$X_x(c)/x\to A(c)$ at each finite $c$; at $c=\infty$, $X/x=1-Y_x/(tx)\to1$. A finite grid on which $A$
varies by at most $\eps$, monotonicity of $X_x(c)$ and a union bound give the uniform statement.
\end{proof}

For a cutoff $h=h(u)$ let $J_u=\set{j:\gamma_j\le h}$, $s_j=(\gamma_j/u)^p$,
$\tau_u^2=\theta\sum_{J_u}s_j^2$, let $Q_u$ be the conditional law of $(tG_j/\gamma_j^p)_{j\in J_u}$ given
$X\le x$, and put $P_u=\bigotimes_{J_u}\mathrm{Gamma}(\theta,1+s_j)$ and
$R_u=\bigotimes_{J_u}\mathrm{Gamma}(\theta,1)$ (second parameter the rate).

\begin{theorem}[Sharp growing-block law; \stU]\label{ana:thm:block}
Let $h=o(u)$.
\begin{enumerate}[label=\textup{(\roman*)}]
\item $d_{\mathrm{TV}}(Q_u,P_u)\to0$; more strongly the joint conditional law of the block and $Y_x$ is
 $o(1)$-close to $P_u\otimes\mathrm{Exp}(1)$.
\item $d_{\mathrm{TV}}(Q_u,R_u)=2\Phi(\tau_u/2)-1+o(1)$. So the block is asymptotically independent
 unit-rate gamma if and only if $\tau_u\to0$, and the distance tends to $1$ if $\tau_u\to\infty$; the
 same holds with the slack included and compared with an independent $\mathrm{Exp}(1)$.
\item If $h\to\infty$, $\tau_u^2\sim\frac{\theta\alpha}{\alpha+2p}N(h)(h/u)^{2p}$.
\end{enumerate}
\end{theorem}

\begin{proof}
Under $\anaPr_t$ write $S=S_B+S_R$ for the block and the independent remainder, with means
$\mu_B+\mu_R=tx$. The block has law $P_u$, and $\mathrm{Var}_t(S_B)\le\theta N(h)=o(\varsigma^2)$. By
\cref{ana:lem:band} the density $g_u$ of $(S_R-\mu_R)/\varsigma$ tends uniformly to $\varphi_1$ and is
uniformly bounded. Relative to $P_u\otimes\mathrm{Exp}(1)$ the conditional law of the block values $r$
and the slack $y$ has density
\begin{equation}\label{ana:eq:blockratio}
 L_u(r,y)=\frac1{\varsigma D_u}\,g_u\Bigl(-\frac{\sum_{J_u}r_j-\mu_B+y}\varsigma\Bigr),
\end{equation}
since the remainder must equal $tx-\sum r_j-y$. Under the reference measure the argument tends to $0$ in
probability, so $L_u\to1$ in probability by \eqref{ana:eq:Du}; $L_u$ is uniformly bounded, so
$\anaEx\abs{L_u-1}\to0$, proving (i).

For (ii), under $R_u$ with coordinates $Z_j$ the log-likelihood of $P_u$ is
$\ell_u=\sum_{J_u}\{\theta\log(1+s_j)-s_jZ_j\}$, with $\max s_j\le(h/u)^p\to0$, mean
$-\tau_u^2/2+O(\max s_j\,\tau_u^2)$ and variance $\tau_u^2$. If $\tau_u\to\tau\in(0,\infty)$, Lyapunov's
condition $\sum\anaEx\abs{s_j(Z_j-\theta)}^3\le C_\theta\max s_j\sum s_j^2\to0$ gives
$\ell_u\Rightarrow\mathcal N(-\tau^2/2,\tau^2)$ (a point mass if $\tau=0$), and
$d_{\mathrm{TV}}(P_u,R_u)=\anaEx_{R_u}(1-e^{\ell_u})_+\to\Phi(\tau/2)-\Phi(-\tau/2)$. If $\tau_u\to\infty$,
the Hellinger affinity $\prod_{J_u}(2\sqrt{1+s_j}/(2+s_j))^\theta=\exp(-\tau_u^2(1+o(1))/8)$ tends to $0$,
and bounds $1-d_{\mathrm{TV}}$. Passing to subsequences in $[0,\infty]$ gives the formula in general;
the triangle inequality and (i) give (ii), and tensoring with the same exponential law preserves total
variation. (iii) is \cref{ana:lem:specsums}.
\end{proof}

A single coordinate, at any height, is never much more likely under the conditioning than under the
tilted law.

\begin{lemma}[Uniform marginal domination; \stU]\label{ana:lem:margdom}
Uniformly in the index $j$ and in Borel sets $E\subseteq(0,\infty)$,
\[
 \anaPr\bigl(tG_j/\gamma_j^p\in E\bigm| X\le x\bigr)\le(1+o(1))\,\anaPr_t\bigl(tG_j/\gamma_j^p\in E\bigr).
\]
In particular the factor is at most $2$ for all sufficiently large $u$.
\end{lemma}

\begin{proof}
Apply \eqref{ana:eq:blockratio} to the block $\set j$, with $g_u$ the density of the standardized
complementary sum, and integrate over the slack. The resulting marginal likelihood ratio is at most
$\sup g_u/(\varsigma D_u)$. By \cref{ana:lem:band} with one index deleted, $\sup g_u\to\varphi_1(0)$
uniformly in $j$, and $\varsigma D_u\to\varphi_1(0)$ by \eqref{ana:eq:Du}.
\end{proof}

Put $\mathcal D(\rho)=2\bigl[\Phi\bigl(\sqrt{-\log(1-\rho)/\rho}\bigr)-
\Phi\bigl(\sqrt{-(1-\rho)\log(1-\rho)/\rho}\bigr)\bigr]$ for $0<\rho<1$, with $\mathcal D(0)=0$ and
$\mathcal D(1)=1$.

\begin{theorem}[Independence with the tilted rates; \stU]\label{ana:thm:macro}
For every positive cutoff $h(u)$, $d_{\mathrm{TV}}(Q_u,P_u)=\mathcal D(V(h/u))+o(1)$. Thus
$d_{\mathrm{TV}}(Q_u,P_u)\to0$ if and only if $h=o(u)$; if $h/u\to c\in(0,\infty)$ the limit is
$\mathcal D(V(c))>0$, and if $h/u\to\infty$ it is $1$.
\end{theorem}

\begin{proof}
Let $h/u\to c\in(0,\infty)$ and $\rho=V(c)$. Bracketing $h/u$ between $c\pm\eps$, the sum and band
lemmas apply to the moving cutoff, so under $P_u$ the centred block sum $Z_u=(S_B-\mu_B)/\varsigma$
tends in law to $\mathcal N(0,\rho)$ and $g_u\to\varphi_{1-\rho}$ uniformly. Integrating
\eqref{ana:eq:blockratio} over the slack (whose shift $y/\varsigma$ has mean $1/\varsigma$), the
likelihood of $Q_u$ relative to $P_u$ is, uniformly in the block, $o(1)$-close to
$r(Z_u)=\varphi_{1-\rho}(Z_u)/\varphi_1(0)=(1-\rho)^{-1/2}\exp(-Z_u^2/(2(1-\rho)))$, a bounded continuous
function. Hence $d_{\mathrm{TV}}(Q_u,P_u)\to\frac12\anaEx\abs{r(Z)-1}$ with $Z\sim\mathcal N(0,\rho)$. Now
$r$ times the density of $\mathcal N(0,\rho)$ is that of $\mathcal N(0,\rho(1-\rho))$; the two densities
cross at $\abs z=\sqrt{-(1-\rho)\log(1-\rho)}$, and integrating the difference between the crossings
gives $\mathcal D(\rho)$. The case $h=o(u)$ is \cref{ana:thm:block}. If $h/u\to\infty$, marginalizing to
the cutoff $bu$ gives $\liminf d_{\mathrm{TV}}\ge\mathcal D(V(b))\to1$ as $b\to\infty$. Subsequences in
$[0,\infty]$ give the general formula.
\end{proof}

\begin{corollary}[The four-fifths transition for completed Dirichlet $L$-functions; \stG]\label{ana:cor:fourfifths}
Take $p=2$ and let the spectrum be the positive ordinates, with multiplicity, of a finite product of
$d$ completed degree-one factors: copies of $\xi(\frac12+w)$, primitive real Dirichlet factors with
their central zeros removed, and conjugate pairs
$H_\chi(w)=\Lambda_\chi(\frac12+w)\Lambda_{\bar\chi}(\frac12+w)/w^{2m}$, $m=\ord_{s=1/2}\Lambda_\chi$,
for nonreal primitive $\chi$ (each pair counting $2$). Then $\anaEx e^{-tX}=(H(0)/H(\sqrt t))^\theta$
for the even entire product $H$, $N(T)\sim\frac d{2\pi}T\log T$, $\alpha=1$, $C_1=\pi/2$, $C_2=\pi/4$,
\[
 A(c)=\frac2\pi\arctan c,\qquad V(c)=\frac2\pi\Bigl(\arctan c+\frac c{1+c^2}\Bigr),\qquad
 x\sim\frac{\theta d\log u}{4u},\qquad\sigma_x^2\sim\frac{\theta d\log u}{8u^3},
\]
and $\tau_u^2\sim\theta d\,h^5\log h/(10\pi u^4)$ as $h\to\infty$. With $h_*(u)=(u^4/\log u)^{1/5}$ and
$h/h_*(u)\to c\in[0,\infty]$, the conditional block $(u^2G_j/\gamma_j^2)_{\gamma_j\le h}$ is at
total-variation distance
\[
 2\Phi\Bigl(c^{5/2}\sqrt{\frac{\theta d}{50\pi}}\Bigr)-1+o(1)
\]
from independent unit-rate gamma coordinates (limits $0$ and $1$ at $c=0,\infty$). In terms of
$n=N(u)$ and $k=N(h)$, if $k/n^{4/5}\to\kappa\in[0,\infty]$ the distance tends to
$2\Phi\bigl(\frac12\sqrt{125\theta\kappa^5/256}\bigr)-1$, independently of $d$. Independence with the
tilted rates holds exactly for $k=o(n)$.
\end{corollary}

\begin{proof}
For a nonreal $\chi$ the functional equations make $H_\chi$ even, and on the real axis
$H_\chi(0)=\abs{\Lambda_\chi^{(m)}(\frac12)/m!}^2>0$. By \cref{src:thm:dgrh} the paired Hadamard product
gives $H_\chi(w)/H_\chi(0)=\prod_{\gamma_j>0}(1+w^2/\gamma_j^2)$, with the positive zeros of both completed
functions; evenness removes the exponential factor. The real factors are as in the proof of
\cref{ana:thm:smalldev}. The count is the sum of the zero counts, and the constants follow from
\cref{ana:thm:condlaw,ana:lem:specsums}. At $h=ch_*(u)$, $\tau_u^2\to2\theta dc^5/(25\pi)$, which gives
the distance by \cref{ana:thm:block}. If $h/h_*\to\infty$, restrict to $h'=\min(h,u/\log u)$, which is
$o(u)$ with $h'/h_*\to\infty$, and marginalize. For the count form,
$N(h)/N(u)^{4/5}\to\frac{4c}5(d/2\pi)^{1/5}$ at $h=ch_*$; substituting $c=\frac{5\kappa}4(2\pi/d)^{1/5}$
gives $\tau^2=125\theta\kappa^5/256$, and monotone bracketing handles general sequences. Regular
variation gives $N(h)/N(u)\to0$ if and only if $h/u\to0$, so \cref{ana:thm:macro} gives the last
claim.
\end{proof}

For example, for $\zeta$ with $\theta=1$ the limiting distances at $c=1,2,3,4$ are about $0.0636$,
$0.3483$, $0.7864$ and $0.9893$. The transition is set by the accumulated squared rate corrections,
not by the largest one, and it involves of order $n^{4/5}$ coordinates. The conductors, root numbers,
multiplicities and individual ordinates affect only the exact saddle and finite-scale corrections.
Monte Carlo comparisons at $u=10^3,10^4$ on a synthetic spectrum with the Riemann--von Mangoldt count
approach the limit slowly, and no finite-$u$ error bound is claimed; the data are in
\begin{center}
\artifact{dirichlet_spectral_conditioning_2026-11-27/data/verification.json}.
\end{center}

For the height carrying half the total, take $\zeta$ ($q=1$) or a primitive real character of conductor
$q$, $p=2$, with central zeros removed and every zero counted with multiplicity, and recall
$N(T)=\frac T{2\pi}(\log\frac{qT}{2\pi}-1)+O_\chi(\log(T+2))$
\cite{BennettMartinOBryantRechnitzer2021Counting}. With the exact saddle $u$ put
\[
 \ell=\log\frac{qu}{2\pi},\qquad J(c)=\int_0^c\frac{\log r}{1+r^2}\,\dd r,\qquad
 M_u(c)=\frac\theta{2\pi u}\bigl(\ell\arctan c+J(c)\bigr),
\]
so that $J(1)=-G$ with Catalan's constant $G$ and $J(\infty)=0$.

\begin{lemma}[Centring on the fluctuation scale; \stG]\label{ana:lem:centring}
Uniformly for $c$ in compact subsets of $(0,\infty)$,
\begin{gather*}
 m_x(c)=M_u(c)+O_\chi\Bigl(\frac{\log u}{u^2}\Bigr)=M_u(c)+o(\sigma_x),\\
 x=\frac{\theta\ell}{4u}+O_\chi\Bigl(\frac{\log u}{u^2}\Bigr),\qquad
 \sigma_x^2=\frac{\theta(\ell-1)}{8u^3}+O_\chi\Bigl(\frac{\log u}{u^4}\Bigr).
\end{gather*}
\end{lemma}

\begin{proof}
Let $R(T)$ be the error of the count, extended to $T\ge0$ with $R(0)=0$. Stieltjes integration by
parts with $f(T)=(T^2+u^2)^{-1}$ gives $m_x(c)=M_u(c)+\theta f(cu)R(cu)-\theta\int_0^{cu}f'(T)R(T)\,\dd T$,
with both error terms $O_\chi(\log u/u^2)$; at $c=\infty$ the boundary term vanishes and $J(\infty)=0$
by $r\mapsto1/r$. For $(T^2+u^2)^{-2}$ the main integrals are $\int_0^\infty(1+r^2)^{-2}\dd r=\pi/4$ and
$\int_0^\infty\log r\,(1+r^2)^{-2}\dd r=-\pi/4$ (write $(1+r^2)^{-2}=\frac12(r/(1+r^2))'+\frac12(1+r^2)^{-1}$
and integrate by parts), with error $O_\chi(\log u/u^4)$. Finally $(\log u/u^2)/\sigma_x\to0$.
\end{proof}

\begin{remark}[A leading-order centring is insufficient]\label{ana:rem:centring}
$m_x(1)-x/2=-\theta G/(2\pi u)+O_\chi(\log u/u^2)$, which divided by $\sigma_x$ grows like
$\sqrt{u/\log u}$. So the centring $xA(c)$, valid on the scale $x$ in \cref{ana:thm:condlaw}(iii), is
not valid in (ii).
\end{remark}

Let $H_x=\min\{\gamma_j:\sum_{\gamma_i\le\gamma_j}G_i/\gamma_i^2\ge X/2\}$, and for $\ell>0$ let $c_\ell>1$
be the unique root of $\ell\arctan c_\ell+J(c_\ell)=\pi\ell/4$; the left side minus the right is $-G$ at
$1$, tends to $\pi\ell/4$ at infinity, and has derivative $(\ell+\log c)/(1+c^2)>0$ for $c\ge1$.

\begin{theorem}[The height carrying half the total; \stG]\label{ana:thm:medianheight}
Conditionally on $X\le x$, as $x\downarrow0$,
\[
 \frac x{\pi\sigma_x}\Bigl(\frac{H_x}u-c_\ell\Bigr)\Longrightarrow\mathcal N\Bigl(0,\frac14-\frac1{\pi^2}\Bigr),
\]
with $c_\ell=1+2G/\ell+2G^2/\ell^2+O(\ell^{-3})$ and $u\pi\sigma_x/x\sim\pi\sqrt{2u/(\theta\ell)}$. The
implicit centre $c_\ell$ must be used: the error of the truncation $1+2G/\ell$, after normalization,
grows like $\sqrt u/\ell^{3/2}$.
\end{theorem}

\begin{proof}
If $c_u\to1$, splitting at $c_uu$ gives variance fractions $v_1=V(1)=\frac12+\frac1\pi$ and $1-v_1$ by
bracketing; the intervals $[u/4,u/2]$ and $[2u,3u]$ lie in the two blocks and supply the Fourier
envelope of \cref{ana:lem:band}. The two-band density argument of the proof of \cref{ana:thm:condlaw}
gives $\bigl((X_x(c_u)-m_x(c_u))/\sigma_x,u^2(x-X)\bigr)\Rightarrow(Z,E)$ with
$Z\sim\mathcal N(0,v_1(1-v_1))$ and $E$ exponential, independent. Let $d_u=\pi\sigma_x/x\to0$, fix $z$, and
take $c_u=c_\ell+zd_u$. By \cref{ana:lem:centring} and the definition of $c_\ell$,
$m_x(c_\ell)-x/2=o(\sigma_x)$, and $M_u'(c)=\frac\theta{2\pi u}\frac{\ell+\log c}{1+c^2}\sim x/\pi$ near
$c=1$, so $(m_x(c_u)-x/2)/\sigma_x\to z$. Also $(x-X)/\sigma_x\to0$ in conditional probability, hence
$(X_x(c_u)-X/2)/\sigma_x\Rightarrow Z+z$. The exact equivalence $H_x\le uc_u\iff X_x(c_u)\ge X/2$ gives the
limit $\anaPr(Z\ge-z)$, and $v_1(1-v_1)=\frac14-\frac1{\pi^2}$. \Cref{ana:lem:centring} gives the height
scale, and Taylor expansion at $c=1$ with $J(1)=-G$, $J'(1)=0$ gives the expansion of $c_\ell$.
\end{proof}

\paragraph{The largest contributions.}
Return to a general regularly varying spectrum, with $\alpha<p$ and $\theta>0$ as before. The
statements above concern fixed coordinates and macroscopic bands; the largest individual terms
\[
 \varrho_j=\frac{tG_j}{\gamma_j^p}=a_j\widetilde G_j
\]
live at a different height. Under $\anaPr_t$ the $\varrho_j$ are independent gamma variables of shape
$\theta$ and rate $1+(\gamma_j/u)^p$. Let $\bar F_\theta(a)=\Gamma(\theta)^{-1}\int_a^\infty s^{\theta-1}e^{-s}\,\dd s$
be the gamma survival function, and define $b_u>0$ and $v_u$ by
\begin{equation}\label{ana:eq:extcentre}
 \mathcal I_u(b):=\sum_j\bar F_\theta\bigl(b[1+(\gamma_j/u)^p]\bigr)=\sum_j\anaPr_t(\varrho_j>b),\qquad
 \mathcal I_u(b_u)=1,\qquad v_u=\frac u{b_u^{1/p}} .
\end{equation}
The series converges uniformly on compact subsets of $(0,\infty)$, and $\mathcal I_u$ is continuous and
strictly decreasing from $\infty$ at $0$ to $0$ at $\infty$, so $b_u$ is well defined. Let
$\nu_{\alpha,p}$ be the probability law on $(0,\infty)$ with density
\[
 \nu_{\alpha,p}(\dd y)=\frac p{\Gamma(\alpha/p)}\,y^{\alpha-1}e^{-y^p}\,\dd y ,
\]
so that $Z\sim\nu_{\alpha,p}$ if and only if $Z^p\sim\mathrm{Gamma}(\alpha/p,1)$. Since
$\sum_j\varrho_j<\infty$ the maximum is attained, and it is unique almost surely, also after
conditioning; order the terms as $\varrho_{I_1}>\varrho_{I_2}>\cdots$ and put
\[
 \Pi_x=\sum_j\delta_{(\varrho_j-b_u,\ \gamma_j/v_u)} .
\]

\begin{theorem}[Conditional marked extremes; \stU]\label{ana:thm:extremes}
\begin{enumerate}[label=\textup{(\roman*)}]
\item $b_u\sim\log N(u)$, and $\dfrac{\Gamma(1+\alpha/p)}{\Gamma(\theta)}N(v_u)\,b_u^{\theta-1}e^{-b_u}\to1$.
\item Conditionally on $X\le x$, $\Pi_x$ converges to a Poisson point process $\Pi$ on
 $\R\times(0,\infty)$ of intensity $e^{-z}\dd z\,\nu_{\alpha,p}(\dd y)$; the convergence includes the
 finite configuration of atoms above any fixed level in the first coordinate. Jointly $Y_x\Rightarrow E_0$,
 exponential of rate $1$ and independent of $\Pi$.
\item For every fixed $k$, conditionally on $X\le x$,
\[
 \Bigl(\bigl(\varrho_{I_i}-b_u,\ (\gamma_{I_i}/v_u)^p\bigr)_{i\le k},\ Y_x\Bigr)\Longrightarrow
 \Bigl(\bigl(-\log(E_1+\dots+E_i),\ H_i\bigr)_{i\le k},\ E_0\Bigr),
\]
 where $E_0,E_1,\dots$ are independent exponential variables of rate $1$, $H_1,H_2,\dots$ are independent
 $\mathrm{Gamma}(\alpha/p,1)$ variables, and the two families are independent. In particular
 $\anaPr(\max_j\varrho_j-b_u\le z\mid X\le x)\to\exp(-e^{-z})$, and the limiting height of the maximizer
 is independent of this Gumbel fluctuation and of the slack.
\end{enumerate}
\end{theorem}

The location law does not depend on the gamma shape $\theta$. The marks converge only in law: their
finite-$u$ distributions are discrete. As before, all randomness is in the weights $G_j$.

\begin{lemma}[Centre and marked intensities; \stU]\label{ana:lem:extintensity}
Part (i) of \cref{ana:thm:extremes} holds. For every $z\in\R$ and every interval $E\subseteq(0,\infty)$,
whose endpoints may be $0$ or $\infty$,
\[
 \sum_{\gamma_j/v_u\in E}\anaPr_t(\varrho_j>b_u+z)\to e^{-z}\nu_{\alpha,p}(E),\qquad
 \max_j\anaPr_t(\varrho_j>b_u+z)\to0 .
\]
\end{lemma}

\begin{proof}
For fixed $b>0$ the first $N(u)$ terms give $\mathcal I_u(b)\ge N(u)\bar F_\theta(2b)\to\infty$, so
$b_u\to\infty$. With $\beta=\max(\theta-1,0)$ the elementary gamma tail bound gives, for $b\ge1$,
$\mathcal I_u(b)\le Cb^\beta e^{-b}\sum_j(1+(\gamma_j/u)^p)^\beta e^{-(\gamma_j/u)^p}\le C'N(u)b^\beta e^{-b}$,
by the dyadic counting of \cref{ana:lem:specsums} with an exponentially decaying integrand. At the root
this gives first $b_u=O(\log N(u))$ and then $b_u\le\log N(u)+O(\log\log N(u))$. For fixed $\eps>0$,
$1=\mathcal I_u(b_u)\ge N(\eps u)\bar F_\theta((1+\eps^p)b_u)$, and
$\bar F_\theta(a)\sim a^{\theta-1}e^{-a}/\Gamma(\theta)$, $N(\eps u)\sim\eps^\alpha N(u)$ give
$\liminf b_u/\log N(u)\ge(1+\eps^p)^{-1}$; so $b_u\sim\log N(u)$ and $v_u\to\infty$.

For $y=\gamma_j/v_u$ and $q=b_u+z$ the argument of the tail is exactly
$q[1+(\gamma_j/u)^p]=b_u+z+(1+z/b_u)y^p$. Uniformly for $y$ in bounded intervals,
\[
 \frac{\bar F_\theta(b_u+z+(1+z/b_u)y^p)}{b_u^{\theta-1}e^{-b_u}/\Gamma(\theta)}\to e^{-z-y^p},
\]
and for fixed $z$ and large $u$ the quotient is at most $C_z(1+y^p)^{\max(\theta-1,0)}e^{-y^p/2}$ on the
whole half-line, by the tail bound for arguments at least $b_u/2$ (for $\theta<1$ the power factor is
bounded since the argument divided by $b_u$ is at least $1/2$). The measures
$N(v_u)^{-1}\sum_j\delta_{\gamma_j/v_u}$ converge to $\alpha y^{\alpha-1}\dd y$, with mass below $\eps$
tending to $\eps^\alpha$, and the exponential envelope with a dyadic regular-variation bound gives
uniformly small tails. Hence
\[
 \mathcal I_u(b_u+z)\sim\frac{N(v_u)b_u^{\theta-1}e^{-b_u}}{\Gamma(\theta)}e^{-z}
 \int_0^\infty\alpha y^{\alpha-1}e^{-y^p}\,\dd y
 =\frac{\Gamma(1+\alpha/p)}{\Gamma(\theta)}N(v_u)b_u^{\theta-1}e^{-b_u}e^{-z},
\]
and the same with the integral restricted to $E$. At $z=0$ the left side is $1$, which is the second
assertion of (i); division gives the marked limit. Each exceedance probability is at most
$\bar F_\theta(b_u+z)\to0$.
\end{proof}

\begin{proof}[Proof of \cref{ana:thm:extremes}]
Under $\anaPr_t$ the coordinates are independent. The joint generating function of the counts in
disjoint rectangles is a product of factors $1+\sum_\ell p_{j,\ell}(w_\ell-1)$; its logarithm differs
from the linear terms by at most $C(\max_j\sum_\ell p_{j,\ell})\sum_{j,\ell}p_{j,\ell}=o(1)$, and the
linear terms converge to the Poisson intensities by \cref{ana:lem:extintensity}. The same argument with
$1-e^{-f}$ gives convergence of Laplace functionals for nonnegative continuous $f$ of compact support.

Fix $R<\infty$. The atoms with mark at most $R$ are determined by the block $\gamma_j\le Rv_u$, and
$Rv_u/u=R/b_u^{1/p}\to0$, so \cref{ana:thm:block}(i) transfers their tilted law to the conditional law,
jointly with an independent exponential slack. By \cref{ana:lem:margdom}, a union bound and
\cref{ana:lem:extintensity}, for every fixed $z_0$,
\[
 \limsup_{u\to\infty}\anaPr\bigl(\exists j:\varrho_j-b_u>z_0,\ \gamma_j/v_u>R\bigm| X\le x\bigr)
 \le2e^{-z_0}\nu_{\alpha,p}((R,\infty)),
\]
which tends to $0$ as $R\to\infty$; marks below $\eps$ are excluded in the same way as $\eps\downarrow0$.
So the truncated limits extend to the full marked process and keep their independence from the slack.
The same bound gives $\limsup\anaPr(\max_j\varrho_j-b_u>A\mid X\le x)\le2e^{-A}$; and restricting marks to
a compact interval $I$ with $\nu_{\alpha,p}(I)>0$, the number of atoms above $-A$ tends to a Poisson
variable of mean $e^A\nu_{\alpha,p}(I)$, which is less than $k$ with probability tending to $0$ as
$A\to\infty$. This justifies ordering the top $k$ atoms. Under $w=e^{-z}$ the unmarked intensity becomes
Lebesgue measure on $(0,\infty)$, whose ordered points are $E_1,E_1+E_2,\dots$; the product intensity makes
the marks independent with law $\nu_{\alpha,p}$ and independent of the arrival times, and taking $p$th
powers of the marks gives (iii).
\end{proof}

For arbitrary regular variation $N(u/b^{1/p})\sim N(u)b^{-\alpha/p}$ may fail, because the multiplier
$b^{-1/p}$ tends to $0$; the implicit centre \eqref{ana:eq:extcentre} avoids this. An explicit centre
uses a reduced counting scale.

\begin{corollary}[Centre from a reduced counting scale; \stU]\label{ana:cor:extcentre}
Let $L_u=\log N(u)$ and $n_u^*=N(u/L_u^{1/p})$. Then
\[
 b_u=\log n_u^*+(\theta-1)\log\log n_u^*+\log\frac{\Gamma(1+\alpha/p)}{\Gamma(\theta)}+o(1).
\]
If moreover $N(u/L_u^{1/p})\sim N(u)L_u^{-\alpha/p}$, for instance if $N(T)\sim cT^\alpha(\log T)^\beta$ with
fixed $c>0$ and $\beta\in\R$, then
\begin{equation}\label{ana:eq:elemcentre}
 b_u=\log N(u)+\Bigl(\theta-1-\frac\alpha p\Bigr)\log\log N(u)+\log\frac{\Gamma(1+\alpha/p)}{\Gamma(\theta)}+o(1).
\end{equation}
\end{corollary}

\begin{proof}
Since $b_u/L_u\to1$, the uniform convergence theorem for regularly varying functions
\cite{BinghamGoldieTeugels1987Regular} gives $N(v_u)/n_u^*\to1$; also $\log n_u^*/\log N(u)\to1$. Taking
logarithms in \cref{ana:thm:extremes}(i) gives the first formula. The additional hypothesis gives
$\log n_u^*=\log N(u)-(\alpha/p)\log\log N(u)+o(1)$.
\end{proof}

\begin{remark}[The reduced scale is needed]\label{ana:rem:slowcount}
Fix $c\ne0$. The integer part of an eventually increasing function
$N(T)\sim T^\alpha\exp(c\log T/\log\log T)$ is the counting function of a locally finite positive spectrum,
regularly varying of index $\alpha$. With $a=\log u$ and $\delta=p^{-1}\log L_u=p^{-1}\log a+O(1)$ one has
$(a-\delta)/\log(a-\delta)-a/\log a\to-1/p$, so $N(u/L_u^{1/p})/(N(u)L_u^{-\alpha/p})\to e^{-c/p}$, and the
exact centre is the right side of \eqref{ana:eq:elemcentre} minus $c/p$, up to $o(1)$. Omitting the
reduced scale can thus move the Gumbel location by a nonzero constant.
\end{remark}

\begin{corollary}[The largest terms come from a band of vanishing mass; \stU]\label{ana:cor:vanishingmass}
Let $v_u=o(h(u))$ and $h(u)=o(u)$. For every fixed $k$, conditionally on $X\le x$,
\[
 \anaPr\bigl(\gamma_{I_1},\dots,\gamma_{I_k}\le h(u)\bigm| X\le x\bigr)\to1,\qquad
 \frac1X\sum_{\gamma_j\le h(u)}\frac{G_j}{\gamma_j^p}\to0,\qquad
 \frac{\max_jG_j/\gamma_j^p}X\sim\frac{b_u}{\theta C_1N(u)}
\]
in conditional probability.
\end{corollary}

\begin{proof}
Tightness of the top $k$ marks and $h/v_u\to\infty$ give the first claim. By \cref{ana:lem:margdom}
the conditional mean of $\sum_{\gamma_j\le h}\varrho_j$ is at most $2\theta N(h)=o(N(u))$, while
$tX=tx-Y_x$ with $tx\sim\theta C_1N(u)$ and $Y_x$ tight; Markov's inequality gives the second claim. The
Gumbel limit gives $\max_j\varrho_j/b_u\to1$.
\end{proof}

\paragraph{The whole bulk path.}
The variance clock $V$ is continuous and strictly increasing from $0$ to $1$. With $V^{-1}(0)=0$ and
$V^{-1}(1)=\infty$ put
\[
 \mathcal B_x(s)=\frac1\varsigma\sum_{\gamma_j\le uV^{-1}(s)}(\varrho_j-\theta a_j)
 =\frac{X_x(V^{-1}(s))-m_x(V^{-1}(s))}{\sigma_x},\qquad0\le s\le1 .
\]
The series converges absolutely almost surely, so $\mathcal B_x\in D[0,1]$, and
$\mathcal B_x(1)=-Y_x/\varsigma$.

\begin{theorem}[Functional bridge, independent of the extremes; \stU]\label{ana:thm:functional}
Conditionally on $X\le x$,
\[
 (\Pi_x,\mathcal B_x,Y_x)\Longrightarrow(\Pi,B,E_0),
\]
with $\Pi$ as in \cref{ana:thm:extremes}, $B$ a standard Brownian bridge, $E_0$ exponential of rate $1$,
the three independent, and $D[0,1]$ carrying the Skorokhod $J_1$ topology. The top-$k$ convergence of
\cref{ana:thm:extremes}(iii) holds jointly with, and independently of, the bulk path.
\end{theorem}

\begin{proof}
\emph{Finitely many cutoffs.} Fix $0<c_1<\dots<c_m$, let $h=Rv_u=o(u)$ be a lower block, put
$\delta=\varsigma^{-1}\sum_{\gamma_j\le h}(\varrho_j-\theta a_j)$, and split the complement into the bands
$(h,c_1u],(c_1u,c_2u],\dots,(c_mu,\infty)$. By \cref{ana:lem:band} their standardized sums have independent
tilted densities $f_{i,u}\to\varphi_{\lambda_i}$ uniformly and in $L^1$, with $\lambda_i=V(c_i)-V(c_{i-1})$,
$\lambda_{m+1}=1-V(c_m)$, $V(c_0)=0$; removing $h=o(u)$ from the first band changes neither its variance
nor its Fourier envelope. With $w_1=\mathcal B_x(V(c_1))$ and $w_i=\mathcal B_x(V(c_i))-\mathcal B_x(V(c_{i-1}))$,
the conditional density relative to the tilted block law times Lebesgue measure in $(w,y)$ is, as in
\eqref{ana:eq:blockratio},
\[
 \frac{e^{-y}\anaone_{y\ge0}}{\varsigma D_u}f_{1,u}(w_1-\delta)\prod_{i=2}^mf_{i,u}(w_i)\,
 f_{m+1,u}\Bigl(-\frac y\varsigma-\sum_{i\le m}w_i\Bigr).
\]
Under the tilted block law $\delta\to0$ in probability. Replace the last density by its normal limit
using uniform convergence, and the first $m$ in turn using $L^1$ convergence (the last one being bounded);
translation continuity in $L^1$ handles $\delta$, and the bounded derivative of the last normal density
handles $y/\varsigma$. With \eqref{ana:eq:Du} this gives convergence in total variation to the product of
the tilted block law, $\mathrm{Exp}(1)$ and the Gaussian-bridge density
$\varphi_1(0)^{-1}\prod_{i\le m}\varphi_{\lambda_i}(w_i)\varphi_{\lambda_{m+1}}(-\sum_iw_i)$, as in the proof of
\cref{ana:thm:condlaw}. The tilted Poisson limit for the block and the restoration of large marks from the
proof of \cref{ana:thm:extremes} then give
$(\Pi_x,(\mathcal B_x(V(c_i)))_i,Y_x)\Rightarrow(\Pi,(B(V(c_i)))_i,E_0)$, with independent components.

\emph{The tilted functional limit.} Under $\anaPr_t$, $\mathcal B_x$ has independent centred increments
and cumulative variance $V_u(s)=\theta\varsigma^{-2}\sum_{\gamma_j\le uV^{-1}(s)}a_j^2$. By
\cref{ana:lem:specsums}, $V_u(s)\to s$ for every $s$, and monotonicity gives
$\eta_u=\sup_s\abs{V_u(s)-s}\to0$. Group equal ordinates: if $\lambda$ occurs $m_\lambda$ times, its jump
is $(a_\lambda/\varsigma)(\mathrm{Gamma}(m_\lambda\theta,1)-m_\lambda\theta)$ with variance
$v_\lambda=m_\lambda\theta a_\lambda^2/\varsigma^2\le2\eta_u$. The fourth central moment of a unit-rate
gamma variable of shape $k$ is $3k^2+6k$, so the fourth moments of the jumps sum to
$3\sum_\lambda v_\lambda^2+6\theta\varsigma^{-4}\sum_ja_j^4\le3\max_\lambda v_\lambda+6\varsigma^{-2}\to0$;
the Lindeberg condition holds even with multiplicities varying with height, and the triangular-array
functional central limit theorem \cite{Whitt2007Proofs} gives a Brownian motion. Directly: partition
$[0,1]$ into $m$ equal intervals; the increment on each has variance at most $m^{-1}+2\eta_u$, and the
fourth cumulants sum to at most $6\varsigma^{-2}$. The $L^4$ martingale maximal inequality on each interval
gives, for the modulus $\omega(f,\delta)=\sup_{\abs{s-s'}\le\delta}\abs{f(s)-f(s')}$,
\[
 \anaPr_t\bigl(\omega(\mathcal B_x,m^{-1})>\eps\bigr)\le C\eps^{-4}\bigl(m^{-1}+2\eta_u+\varsigma^{-2}\bigr);
\]
near $s=1$ truncate the row and pass to the $L^4$ limit.

\emph{Transfer on initial intervals.} Fix $r<1$. The coordinates with $\gamma_j\le uV^{-1}(r)$ form a
prefix; the standardized complementary sum has a density with bounded supremum (keep a fixed band beyond
$V^{-1}(r)$ and use the Fourier envelope of \cref{ana:lem:band}; convolution does not increase a
supremum). The marginal form of \eqref{ana:eq:blockratio} and $\varsigma D_u\to\varphi_1(0)>0$ give
$\anaPr(E\mid X\le x)\le C_r\anaPr_t(E)$ for all events $E$ of the prefix and all large $u$, so tilted
tightness on $[0,r]$ transfers.

\emph{The endpoint.} For $r\ge1/2$ the complement of the tail $\gamma_j>uV^{-1}(r)$ contains the fixed band
$(V^{-1}(1/4),V^{-1}(1/2)]$, so the same argument gives a constant $C_0$, independent of $r\ge1/2$ and of
large $u$, with $\anaPr(E\mid X\le x)\le C_0\anaPr_t(E)$ for events of the tail. The $L^2$ maximal
inequality for the reversed tail (truncate, then pass to the limit) gives
\[
 \anaPr\Bigl(\sup_{s\ge r}\abs{\mathcal B_x(1)-\mathcal B_x(s)}>\eps\Bigm| X\le x\Bigr)
 \le\frac{4C_0}{\eps^2}\bigl(1-r+\eta_u\bigr),
\]
which is small for $u\to\infty$ and then $r\uparrow1$. With the modulus on $[0,r]$ (insert the value at $r$
for pairs straddling it) this proves tightness on $[0,1]$ with continuous subsequential limits, and
$\mathcal B_x(1)=-Y_x/\varsigma\to0$.

\emph{Identification.} Marginal tightness gives joint tightness; evaluation at fixed times is continuous
at continuous limits, and the finite-dimensional limits of the first step, with the endpoint values $0$,
identify every subsequential limit as $(\Pi,B,E_0)$ with independent components. The tail bounds of the
proof of \cref{ana:thm:extremes} allow the top-$k$ map jointly.
\end{proof}

\paragraph{Dirichlet spectra.}
Take the spectrum of \cref{ana:cor:fourfifths}: the positive ordinates, with multiplicity, of an even
product $H$ of $d$ completed degree-one factors with central zeros removed, so that by
\cref{src:thm:dgrh} $\anaEx e^{-tX}=(H(0)/H(\sqrt t))^\theta$ for $X=\sum_jG_j/\gamma_j^2$, and
$p=2$, $\alpha=1$, $N(T)\sim\frac d{2\pi}T\log T$ \cite{BennettMartinOBryantRechnitzer2021Counting}.

\begin{corollary}[Half-normal height of the largest Dirichlet contribution; \stG]\label{ana:cor:dirextremes}
With the exact saddle $u=\sqrt t$, the centre
\[
 b_u=\log u+\Bigl(\theta-\frac12\Bigr)\log\log u+\log\frac d{4\sqrt\pi\,\Gamma(\theta)}+o(1)
\]
is valid, and conditionally on $X\le x$
\[
 u^2\max_j\frac{G_j}{\gamma_j^2}-b_u\Longrightarrow G,\qquad
 \frac{\sqrt{b_u}\,\gamma_{I_1}}u\Longrightarrow Z,
\]
where $\anaPr(G\le z)=e^{-e^{-z}}$ and $\anaPr(Z\in\dd y)=\frac2{\sqrt\pi}e^{-y^2}\anaone_{y>0}\,\dd y$,
jointly with $u^2(x-X)\Rightarrow E_0$ and with the bridge of \cref{ana:thm:functional} (for which
$V(c)=\frac2\pi(\arctan c+\frac c{1+c^2})$), all four limits independent. Thus $Z$ is the absolute value of
a centred normal variable of variance $1/2$, and
\[
 x\sim\frac{\theta d\log u}{4u},\qquad\frac{\max_jG_j/\gamma_j^2}X\sim\frac4{\theta d\,u}
\]
in conditional probability. For $\zeta$ with exponential weights ($d=\theta=1$) and
$b_0(u)=\log u+\frac12\log\log u-\log(4\sqrt\pi)$, this reads
$\bigl(u^2\max_jE_j/\gamma_j^2-b_0(u),\ \sqrt{\log u}\,\gamma_{I_1}/u,\ u^2(x-X)\bigr)\Rightarrow(G,Z,E_0)$.
\end{corollary}

\begin{proof}
Here $C_1=\pi/2$ and $C_2=\pi/4$. Substituting $N(u)\sim du\log u/(2\pi)$ into \eqref{ana:eq:elemcentre},
with $\Gamma(3/2)=\sqrt\pi/2$, gives the centre, and $\nu_{1,2}$ is the displayed half-normal law. The
saddle asymptotic is \cref{ana:cor:fourfifths}, and the ratio follows from
\cref{ana:cor:vanishingmass}. Since $b_u\sim\log u$, $\sqrt{b_u}$ may be replaced by $\sqrt{\log u}$.
\end{proof}

So about half the mass lies below height $u$ (the profile is $A(c)=\frac2\pi\arctan c$), while the largest
term sits at height of order $u/\sqrt{\log u}$. With $h=u/(\log u)^{1/4}$, \cref{ana:cor:vanishingmass}
shows that every fixed number of largest terms lies below $h$ with probability tending to $1$, although
the mass below $h$ is a vanishing fraction of $X$.

The elementary centre converges slowly. For exponential weights a two-term count gives a much sharper
implicit centre, with an explicit enclosure; this concerns the deterministic normalization only and
gives no rate for the conditional limit.

\begin{proposition}[Heat-trace centring; \stU]\label{ana:prop:refinedcentre}
Let $p=2$, $\theta=1$, and $N(T)=c_NT(\log T+a_N-1)+O(\log(T+2))$ with fixed $c_N>0$ and $a_N$. Put
$k_N=c_N\sqrt\pi/2$ and $\kappa_N=a_N-\gamma_{\mathrm E}/2-\log2$, and let $B_u$ be the unique solution in
$0<B<u^2e^{2\kappa_N}$ of
\begin{equation}\label{ana:eq:refinedroot}
 e^B\sqrt B=k_Nu\Bigl[\log\frac u{\sqrt B}+\kappa_N\Bigr].
\end{equation}
Then $B_u\sim\log u$ and $b_u=B_u+O(\sqrt{\log u}/u)$. Explicitly, let $R(T)=N(T)-c_NT(\log T+a_N-1)$,
$R(0)=0$. If $\abs{R(T)}\le C\log(T+2)$ for all $T\ge0$, put $v_B=u/\sqrt{B_u}$ and
$\eps_u=Ce^{-B_u}[\log(v_B+2)+\sqrt\pi/2]$; whenever $\eps_u<1$,
\[
 B_u+\log(1-\eps_u)\le b_u\le B_u+\log(1+\eps_u).
\]
If instead $\abs{R(T)}\le C_0$ for all $T$, the same holds with $\eps_u=C_0e^{-B_u}$, and
$b_u-B_u=O(1/(u\sqrt{\log u}))$.
\end{proposition}

\begin{proof}
Let $\mathcal H(v)=\sum_je^{-(\gamma_j/v)^2}$. Stieltjes integration gives
\[
 \mathcal H(v)=c_N\int_0^\infty e^{-(T/v)^2}(\log T+a_N)\,\dd T+\frac2{v^2}\int_0^\infty Te^{-(T/v)^2}R(T)\,\dd T
 =k_Nv(\log v+\kappa_N)+\mathcal R(v),
\]
the boundary terms vanishing since $R(0)=0$; here $\int_0^\infty e^{-y^2}\log y\,\dd y=\frac{\sqrt\pi}2(-\frac{\gamma_{\mathrm E}}2-\log2)$,
by differentiating $\int_0^\infty y^{s-1}e^{-y^2}\dd y=\Gamma(s/2)/2$ at $s=1$ with
$\psi(1/2)=-\gamma_{\mathrm E}-2\log2$. Since $vy+2\le(v+2)(1+y)$ and $\log(1+y)\le y$,
$\abs{\mathcal R(v)}\le2C\int_0^\infty ye^{-y^2}\log(vy+2)\,\dd y\le C[\log(v+2)+\sqrt\pi/2]$, and
$\abs{\mathcal R(v)}\le C_0$ under a bounded error. (An asymptotic count extends to a global logarithmic
bound by enlarging $C$ on a compact interval.) For $\theta=1$, $\mathcal I_u(q)=e^{-q}\mathcal H(u/\sqrt q)$.
The function $F_u(q)=e^{-q}k_Nuq^{-1/2}(\log u-\frac12\log q+\kappa_N)$ is positive on
$(0,u^2e^{2\kappa_N})$, runs from $\infty$ to $0$, and has
$(\log F_u)'=-1-\frac1{2q}-\frac1{2q(\log u-\frac12\log q+\kappa_N)}<0$, so $F_u(B)=1$, which is
\eqref{ana:eq:refinedroot}, has one root; $F_u(a\log u)\sim k_Na^{-1/2}u^{1-a}\sqrt{\log u}$ gives
$B_u\sim\log u$. At $q=B_u$ one gets $\abs{\mathcal I_u(B_u)-1}\le\eps_u$. Termwise,
\[
 \frac{\dd}{\dd q}\log\mathcal I_u(q)=-\frac{\sum_j[1+(\gamma_j/u)^2]e^{-q[1+(\gamma_j/u)^2]}}{\sum_je^{-q[1+(\gamma_j/u)^2]}}\le-1,
\]
so $0\le b_u-B_u\le\log\mathcal I_u(B_u)$ if $\mathcal I_u(B_u)\ge1$ and $0\le B_u-b_u\le-\log\mathcal I_u(B_u)$
otherwise, which is the enclosure. Finally \eqref{ana:eq:refinedroot} gives
$e^{-B_u}=1/(k_Nv_B(\log v_B+\kappa_N))$, so $\eps_u=O(1/v_B)=O(\sqrt{\log u}/u)$, and
$O(1/(v_B\log v_B))=O(1/(u\sqrt{\log u}))$ under a bounded error.
\end{proof}

Writing $L=\log u$, expansion of \eqref{ana:eq:refinedroot} gives
$B_u=L+\frac12\log L+\log k_N+(\kappa_N-\frac34\log L-\frac12\log k_N)/L+O((\log L)^2/L^2)$, which explains
the slow approach to the elementary centre.

\begin{remark}[Conductor dependence of the centre; \stG]\label{ana:rem:conductorcentre}
For the product of \cref{ana:cor:dirextremes} let $\mathcal Q=\prod_{\ell\le d}q_\ell$ be the product of
the conductors, conjugate factors counted separately and $\zeta$ given conductor $1$. The two-term zero
counts give $c_N=d/(2\pi)$ and $a_N=(\log\mathcal Q)/d-\log(2\pi)$, unchanged by removing finitely many
central zeros, so for exponential weights $B_u$ is the root with positive bracket of
\[
 e^B\sqrt B=\frac{du}{4\sqrt\pi}\Bigl[\log\Bigl(\frac{\mathcal Q^{1/d}u}{4\pi\sqrt B}\Bigr)-\frac{\gamma_{\mathrm E}}2\Bigr],\qquad
 b_u=B_u+O_H\Bigl(\frac{\sqrt{\log u}}u\Bigr).
\]
The estimate is for a fixed product; no uniformity in the conductor or the degree is claimed.
\end{remark}

On the synthetic spectrum $\gamma_j=2\pi j/W(j/e)$ ($W$ the Lambert function), whose count differs from
$\frac T{2\pi}(\log\frac T{2\pi}-1)$ by at most $1$, the numerically computed exact centre at $u=10^5$ is
$10.40117$; it differs from the elementary centre by $0.375$ and from $B_u$ by $1.4\cdot10^{-5}$, inside
the enclosure with $\eps_u=e^{-B_u}$ (floating-point diagnostics, not interval certificates; these are not
zeta zeros); see \artifact{spectral_extremes_2026-11-27/results/refined_center_validation.json}.

Conditioning small Gaussian norms and conditional bridge limits for weighted sums have classical
precedents \cite{VanCampenhoutCover1981Maximum,MayerWolfZeitouni1993Probability,SobolevShklyaev2020Large};
Gumbel limits for maxima of gamma variables are classical \cite{BoseDasguptaMaulik2008Maxima}, and
extremal Fourier intensities with power-law rates, without a small-total constraint, are studied in
\cite{GyorgyiHoldsworthPortelliRacz2003Statistics}.
All randomness here comes from the gamma weights; nothing is asserted about randomness, spacing,
simplicity or independence of the zeros, and no uniformity in a growing conductor or degree is
claimed.

\subsection{A finite certificate for Euler's Gamma product}\label{ana:sec:eulergamma}
{\raggedright\sources{gamma-euler-quartic}}

Archimedean factors such as those in \cref{ana:thm:smalldev} are values of the Gamma function. For
$N\ge1$ and $\Re z>0$ put
\[
 E_N(z)=\frac{N^zN!}{z(z+1)\cdots(z+N)},\qquad A=\frac{z(z+1)}2,\qquad D=-\frac{z(z+1)(2z+1)}{12},\qquad
 C=\frac{A^2}3,
\]
so that $E_N(z)\to\Gamma(z)$ (Euler's product), and define the polynomial multiplier
\[
 P_N=1+\frac AN+\frac{D+A^2/2}{N^2}+\frac{C+AD+A^3/6}{N^3},
\]
the cubic truncation of $\exp(A/N+D/N^2+C/N^3)$; its coefficients are the classical Gamma-ratio
coefficients. The point is an explicit finite error bound for this particular sampler that uses no
logarithm of a Gamma value.

\begin{proposition}[Fourth-order Euler sampler; \stU]\label{ana:prop:eulergamma}
Let $R\ge1$ be an integer, $a>0$, $M=R+1$ and $B=\max\{1/a,(R-1)!\}$. Uniformly for $\abs z\le R$ and
$\Re z\ge a$, and for $N\ge2M^2$,
\[
 \abs{\Gamma(z)-E_N(z)P_N(z)}\le\frac{3BM^8}{N^4},\qquad\abs{P_N(z)}\le\frac32 .
\]
\end{proposition}

\begin{proof}
\emph{Size.} With $x=\Re z$, termwise comparison gives $\abs{E_N(z)}\le E_N(x)$, and $\log E_N(x)$ is
convex in $x>0$, its nonlinear terms being $-\log(x+j)$. At $b=\min(a,1)$, concavity of $t^b$ gives
$\prod_{j\le N}(1+b/j)\ge(N+1)^b$, so $E_N(b)\le1/b$; at the integer $R$, $E_N(R)\le(R-1)!$. Thus $B$
bounds $\abs{E_N(z)}$ and its limit.

\emph{Telescope.} Put $w=z+1$, $t=1/n$ and $S_n=A/n+D/n^2+C/n^3$. The product gives $E_{n+1}=E_ne^{L_n}$
with $L_n=w\log(1+t)-\log(1+wt)$, by the logarithm series for $\abs{wt}\le1/2$. Its coefficients in
degrees $2,3,4$ are $A$, $(w-w^3)/3$, $(w^4-w)/4$, and those of $S_n-S_{n+1}$ are $A$, $-A+2D$,
$A-3D+3C$, which agree. The rest of the logarithm series is at most
$(M+M^5)t^5/[5(1-Mt)]\le M^5t^5$. In $S_n-S_{n+1}$ the tail after degree $4$ is at most
$(\abs A+4\abs D+6\abs C)t^5$, from the geometric series for $(1+t)^{-1}$ and the Taylor bounds for
$(1+t)^{-2},(1+t)^{-3}$, whose next coefficients are $4$ and $6$ in absolute value. As $\abs A\le M^2/2$,
$\abs D\le M^3/6$ and $\abs C\le M^4/12$, the defect $d_n=L_n-S_n+S_{n+1}$ satisfies
$\abs{d_n}\le2M^5/n^5$. For $n\ge4$, $1/n^4-1/(n+1)^4\ge2/n^5$, so every finite sum of defects from $N$
on is at most $M^5/N^4$. The exact recursion $E_K=E_N\exp(S_N-S_K+\sum_{n=N}^{K-1}d_n)$, $S_K\to0$ and
$E_K\to\Gamma(z)$ give $\Gamma=E_Ne^{S_N+d}$ with $\abs d\le M^5/N^4$.

\emph{Rational multiplier.} With $u=M^2/N\le1/2$, the coefficients of $e^{S_N}$ as a series in $u$ are
bounded by those of $\exp(u/2+u^2/6+u^3/12)$, whose sum at $u=1$ is $e^{3/4}<5/2$. Discarding powers
$u^{\ge4}$ costs at most $\frac52u^4$, and the retained cubic is exactly $P_N$; its coefficient bounds
$1/2,7/24,3/16$ give $\abs{P_N}\le3/2$. Since $\abs{S_N}+\abs d\le1$,
$\abs{e^{S_N+d}-e^{S_N}}\le3\abs d$, and $\frac52M^8+3M^5<3M^8$ for $M\ge2$.
\end{proof}

For $R=8$, $a=1/2$ and $N=2^{42}$, an observation of $E_N$ with absolute error at most $2^{-130}$ gives
$\Gamma(z)$ to within $2^{-128}$, since $3\cdot5040\cdot9^8\,2^{-168}+\frac32\,2^{-130}<2^{-128}$; for
the reciprocal $1/E_N$, $156$ absolute bits suffice before exact inversion, as $\abs{1/E_N}\ge1/5040$
and the inverse error is at most $2\cdot5040^2$ times the input error when that is at most
$1/(2\cdot5040)$. The uncorrected value $E_N(1)=N/(N+1)$ would need $N\ge2^{128}-1$ for the same
accuracy. Exact rational checks at $680$ integer inputs and $2232$ nonreal rational local defect
enclosures are recorded in \artifact{gamma-euler-quartic/verification.txt}. The proposition is
independent of \cref{ana:thm:window}, whose Gamma averages are probability measures on the logarithmic
variable rather than values of $\Gamma$.
